% =====================================================================
% Manuscript: Emergent Carrying Capacity, the Biocapacity Ratio, and the Productivity Illusion: Deficit-Driven Collapse in a Delayed Coupled Human–Environment Model
% Author: Amin Abaee (Independent Researcher)
% Compiled with LuaLaTeX (fontspec + unicode-math) - error-free.
% =====================================================================
\documentclass[11pt]{article}

\usepackage[a4paper,margin=1in]{geometry}
\usepackage{fontspec}
\usepackage{amsmath,amssymb}
\usepackage{unicode-math}
\usepackage{booktabs}
\usepackage{array}
\usepackage{tabularx}
\usepackage{graphicx}
\usepackage{enumitem}
\emergencystretch=3em
\setlength{\tolerance}{2000}
\usepackage[colorlinks=true,urlcolor=blue,linkcolor=blue,citecolor=blue]{hyperref}

\setmainfont{DejaVu Serif}
\setmonofont{DejaVu Sans Mono}
\setmathfont{Latin Modern Math}

\hypersetup{
  pdfauthor={Amin Abaee},
  pdftitle={Emergent Carrying Capacity, the Biocapacity Ratio, and the Productivity Illusion},
}

\title{Emergent Carrying Capacity, the Biocapacity Ratio, and the Productivity Illusion: Deficit-Driven Collapse in a Delayed Coupled Human–Environment Model}
\author{
  \textbf{Amin Abaee}\\
  \textit{Independent Researcher}\\
  \href{mailto:amin_abaee@ut.ac.ir}{amin\_abaee@ut.ac.ir}\\
  \href{https://orcid.org/0000-0002-0019-1842}{ORCID: 0000-0002-0019-1842}
}
\date{September 6, 2026}

\begin{document}
\maketitle

\begin{abstract}
Coupled human–environment models typically treat carrying capacity as an imposed, static ceiling. We show that it is instead emergent from the interaction of environmental stock, technological productivity and per-capita demand; treating it as fixed can produce a false sense of environmental health. We develop a minimal deterministic model in which environmental assets (forests, soils, water systems) function as natural capital that yields a regenerative flow. Cumulative overshoot is tracked as ecological debt that erodes that yield; environmental regeneration and demographic response act through two distinct delays. Three results follow. First, under full harvest (\ensuremath{\sigma{}} = 1) the constant-parameter subsystem admits a one-parameter family of equilibria $P = B(A)/e$ rather than an isolated attractor; for the representative parameter set every interior point is monotonically unstable — a positive real eigenvalue for every delay, no imaginary-axis crossing, for $b/b_{G} + G\prime{}(A*) > r$. Collapse onset is therefore a structural vicious cycle, not a delay-ratio Hopf bifurcation, and is not controlled by a single stability index. Two mechanisms operate: debt-driven collapse even in the absence of delay, and a delay-amplified transient that shrinks the basin of attraction. Second, the operating boundary is the biocapacity ratio $R_{B} = 1$ (footprint equal to total biocapacity); the flow- yield ratio $R_{A} = 1$ is a leading but non-causal signal. Once the regeneration lag is long, neither ratio warns. Third, a sufficiently early technology wave can raise measured biocapacity while the stock continues to fall, but only inside a narrow window of about five years that exists solely for small initial deficits and vanishes beyond modest overshoot. This is a productivity illusion, not improving health: yield gains saturate while debt compounds, and the same technology simultaneously raises biocapacity, sustainable population and aggregate demand, tending to increase cumulative debt under overshoot. The decisive lever is the regeneration delay; neither a rising harvest nor an early- sounding ratio reliably signals collapse.
\end{abstract}

\noindent\textbf{Keywords: }* carrying capacity; ecological footprint; biocapacity; time delay; ecological debt; sustainability; delayed feedback

\section{Introduction}

Carrying capacity is one of ecology's most enduring yet elusive ideas. In classical population biology it is a fixed number — the ceiling that caps growth. Imported into global human ecology it is usually treated as a static boundary, the maximum number of people the Earth can support (Cohen, 1995). Such treatments overlook a feedback that matters: if a population overshoots its resource base, it degrades the very environment that sustains it and thereby lowers the future ceiling. The ceiling is a moving target. Imagine an orchard. Each year the trees bear fruit, and this fruit can be gathered without harming the trees themselves. This is a \emph{flow}; ecologists call the annual regenerative flow of productive land and water biocapacity. The trees themselves remain standing — a stock of natural capital — and it is this stock that generates the flow. Taking the fruit at or below the rate at which the trees regrow uses the flow, which is biocapacity; cutting into the standing stock faster than it regenerates is liquidation of capital, which reduces the future flow rather than contributing to it. Suppose the harvesters take more than the orchard can regrow. At first nothing seems wrong — the trees still stand, the harvest still comes — but the surplus is not free. It is paid for by cutting branches, exhausting soil, and eventually felling trees. A harvest that looks stable is masking a shrinking asset. We call this the productivity illusion: the mistaken belief that a steady or rising harvest signals a healthy orchard. It sits at the heart of a long-standing debate in sustainability science between those who hold that technology can always substitute for natural capital and those who argue that substitution has hard limits. The Global Footprint Network turns this bookkeeping into two accounting quantities (Borucke et al., 2013). Biocapacity is the fruit; the ecological footprint is the demand placed on it. When the footprint exceeds biocapacity, an ecological deficit occurs, analogous to drawing down a bank account: the annual balance is a flow, but its cumulative effect erodes the underlying stock of natural capital, which in turn determines future biocapacity. The very quantity that defines the ceiling is eroded by sustained overshoot. Nor do the orchard and its harvesters respond instantly. Forest loss, soil degradation and climate change unfold over decades, while demographic patterns and social norms are slow to change. The system is therefore a delay differential equation, a class widely used to capture feedback delays in ecology (Kuang, 1993), and this is why civilisations can prosper for long periods while quietly undermining their productive base, only to suffer abrupt decline (Tainter, 1988; Diamond, 2005). The orchard also has an everyday analogue. A system can bear overload for a while — an elevator rated for ten people will still manage fourteen, showing little visible strain before the cable gives way — while its underlying integrity is quietly eroding. Judged by what it delivers (biocapacity, the harvest) the elevator still runs; judged by its reserve (the stock, the cable) it is already failing. This gap is the productivity illusion, and the question this paper addresses is not whether such masking can occur but when it can be sustained and when monitoring fails to reveal it. The wider literature frames the same problem. The most prominent global human–environment model is the World3 system of \emph{The Limits to Growth} (Meadows et al., 1972, 2004) and \emph{Dynamics of Growth in a Finite World} (Meadows et al., 1974), which produces overshoot scenarios through complex feedbacks among numerous interacting stocks; a rich literature on coupled human–environment and resource-harvest dynamics has since developed (Brander \& Taylor, 1998; Anderies, 2000). Estimates of Earth's carrying capacity span an enormous range, from well under one billion to over one trillion people (Cohen, 1995; van den Bergh \& Rietveld, 2004), and some scholars argue the very notion of a fixed carrying capacity is inapplicable to humans (Simon, 1996). Meyer (1994) and Meyer \& Ausubel (1999) instead model carrying capacity as itself logistically varying, and Dasgupta (2021) formalises natural capital as a depreciating asset. These diverging perspectives all point to the same need: frameworks in which carrying capacity is not an exogenous input but an emergent property of the system. The model we develop makes the orchard precise. Biocapacity is the sum of a flow yield — separable from the stock, taken without reducing it — and the capital's growth — removable only by liquidating the base. The share of the two, the flow share, governs how the system responds to demand and tells us whether a rise in measured biocapacity is a true recovery or a mask over a shrinking stock. Because the stock is the only way to meet a shortfall, overshoot feeds damage back onto the supply: demand in excess of the flow yield is paid from the base, and cumulative overshoot accumulates as an ecological debt that steadily erodes yield. Two delays complete the picture — the regeneration delay before the orchard recovers, and the demographic delay before a population responds to a changing ceiling. Three conclusions follow, and each is a statement about the orchard. First, a sustainable balance is not one point but a family of balances, and once demand presses past the renewable flow the system begins a self-reinforcing loss that needs no oscillatory mechanism to start: the decline is structural, not a delicate, delay-tuned wobble. Second, the boundary that separates recovery from collapse is when total demand equals total biocapacity; the more commonly watched flow-only ratio is an early but non-decisive signal, and when the regeneration delay is long even that can stay silent. Third, an early enough technology wave can make measured biocapacity rise while the stock falls — but only within a narrow window, about five years, and only while the initial deficit is small; beyond a modest overshoot the window closes. The same technology that lifts biocapacity also lifts the sustainable population and hence total demand, so under overshoot it tends to increase cumulative debt. The message is therefore cautionary, and it names the levers precisely. The lever that matters is the regeneration delay — shorten it, and adopt the protective, effort-reducing institutional stance — while the levers that do not exist include rescue-by-stock and any early-warning signal. Rising biocapacity need not mean improving environmental health; it may be the productivity illusion that precedes a reckoning. The paper remains a conceptual, stylised illustration of the accounting logic, not a calibrated forecast: parameters are representative, the eight predictions it states are falsifiable hypotheses to be tested rather than results to be asserted, and the consequences of the model's deliberately minimal, Allee-free choices are discussed in §Model scope and §Limitations. ---

\section{Model formulation}

\subsection{Variables and units}

\begin{table}[htbp]
\centering
\small
\begin{tabularx}{\textwidth}{>{\raggedright\arraybackslash}X>{\raggedright\arraybackslash}X>{\raggedright\arraybackslash}X>{\raggedright\arraybackslash}X}
\toprule
Symbol & Meaning & Unit & State? \\
\midrule
$A$ & productive biocapacity area (the "trees") & ha & yes \\
$P$ & human population & cap & yes \\
$D$ & accumulated ecological debt (degradation channel) & gha\ensuremath{\cdot{}}yr & yes \\
$b$ & flow yield per unit area (fruit per tree) & gha\ensuremath{\cdot{}}ha⁻¹\ensuremath{\cdot{}}yr⁻¹ & derived \\
$b_{G}$ & value of one hectare of standing stock & gha\ensuremath{\cdot{}}ha⁻¹ & parameter \\
$G(A)$ & regeneration / recruitment rate & ha\ensuremath{\cdot{}}yr⁻¹ & no \\
$B$ & biocapacity = $bA + b_{G} G(A)$ & gha\ensuremath{\cdot{}}yr⁻¹ & no \\
$E$ & footprint = $eP$ & gha\ensuremath{\cdot{}}yr⁻¹ & no \\
$K$ & carrying capacity = \texttt{B/e} & cap & algebraic, not a state \\
$e$ & per-capita footprint & gha\ensuremath{\cdot{}}cap⁻¹\ensuremath{\cdot{}}yr⁻¹ & param \\
$\rho{}$ & regeneration rate & yr⁻¹ & param \\
$r$ & per-capita population growth rate & yr⁻¹ & param \\
$\alpha{}$ & degradation rate & (gha\ensuremath{\cdot{}}yr)⁻¹ & param \\
$\eta{}$ & debt-repayment rate & yr⁻¹ & param \\
$\tau{}_{g}$ & recruitment (regeneration) delay & yr & param \\
$\tau{}_{p}$ & demographic (generation) delay & yr & param \\
$A_{max}$, $b_{0}$, $A_{ext}$ & carrying-area ceiling, base yield, extinction floor & ha, gha\ensuremath{\cdot{}}ha⁻¹\ensuremath{\cdot{}}yr⁻¹, ha & param \\
$T_{b}$ & technology-driven yield step (the yield-enhancing wave $T_{b}(t)$) & gha\ensuremath{\cdot{}}ha⁻¹\ensuremath{\cdot{}}yr⁻¹ & param \\
$\Delta{}b$ & amplitude of the yield wave (asymptotic rise of $T_{b}$) & gha\ensuremath{\cdot{}}ha⁻¹\ensuremath{\cdot{}}yr⁻¹ & param \\
$\kappa{}$ & steepness / rate of the yield wave & yr⁻¹ & param \\
$t_{wave}$ & time of the yield wave's midpoint & yr & param \\
\bottomrule
\end{tabularx}
\end{table}

Because $A$ is in physical hectares it is independent of yield, so $b = B/A$ is definable and the decomposition $B = b\cdot{}A + b_{G}\cdot{}G(A)$ is identifiable. (With a stock measured directly in global hectares — as in the National Footprint Accounts — a hectare already embeds $b$, and the decomposition is not identifiable.) The full symbol table with the units of every quantity is given in the Supplementary Information (§S1.2). \textbf{Units — the flow/stock convention.} $E$ (footprint) and $B$ (biocapacity) are per-annum flows and are written $gha\cdot{}yr^{-1}$; they are the amount of global-hectare-equivalent demand (or supply) in one year. The accumulated ecological debt $D$ is their time-integral over the deficit, $D = \int{}[E - B]_{+} dt$, and therefore carries $gha\cdot{}yr$ (global-hectare-years) — a \emph{stock} of past overshoot, not a rate. The two are never interchanged: a $gha\cdot{}yr^{-1}$ rate enters the rate equations (4) and (7) directly, while the $gha\cdot{}yr$ accumulation $D$ enters only the degradation law (8) as $e^{-\alpha{}D}$. The Global Footprint Network reports the footprint and biocapacity of a given year in global hectares (gha) with the "per year" basis left implicit; the model makes that basis explicit ($gha\cdot{}yr^{-1}$) so that $E$ and $B$ can be differentiated in time. The per-capita footprint $e$ is likewise $gha\cdot{}cap^{-1}\cdot{}yr^{-1}$ (per-person annual demand), so $E = eP$ is a flow.

\subsection{Governing equations}

\begingroup\small\begin{verbatim}
Regeneration:  G(A) = ρ A (1 − A/A_max)                          (1)
Biocapacity:   B    = b A + b_G G(A)                            (2)
Footprint:     E    = e P                                       (3)
Stock:         dA/dt = G(A(t−τ_g)) − [E(t) − σ b A(t)]₊ / b_G   (4)
Population:    dP/dt = r P(t) [1 − P(t) / K(t−τ_p)]              (5)
Carrying cap:  K    = B / e                                     (6)   [algebraic]
Debt:          dD/dt = [E − B]₊ − η D                           (7)
Degradation:   b    = (b₀ + T_b(t)) e^{−α D(t)}                 (8)
Technology:    T_b  = Δb / (1 + e^{−κ(t−t_wave)})               (9)
\end{verbatim}\endgroup

\textbf{Notation and conventions.} Four conventions are used throughout and are stated here so that the equations are unambiguous:

\begin{itemize}
\item \textbf{$\sigma{}$ is a dimensionless fraction}, $0 \leq{} \sigma{} \leq{} 1$, the human-available share of the flow. $\sigma{} = 1$ means the \emph{entire} flow yield $bA$ may be harvested without damaging the base; $\sigma{} < 1$ reserves a share $(1-\sigma{})B$ (a Half-Earth / reservation policy, §9), and $\sigma{}$ is then set together with the population cap so that $E \leq{} \sigma{}\cdot{}B$. The factor $\sigma{} b A(t)$ in (4) is the safe flow harvest; only the excess $[E(t) - \sigma{} b A(t)]_{+}$ draws down the stock. \textbf{Consequence:} the one-parameter equilibrium family $P = B(A)/e$ of §4.3 and §4.5, and the neutral continuum $D(0)=0$ on which it rests, are derived under $\sigma{} = 1$. Under a reservation $\sigma{} < 1$ the equilibrium locus shifts to $P = (\sigma{}bA + b_{G} G(A))/e < B/e$, i.e. reserving a share removes $(1-\sigma{})bA$ from the sustainable population (verified in §4.3 / §S3.3).
\item \textbf{$[x]_{+}$ is the positive part, $[x]_{+} = max(x, 0)$.} In (4) and (7) it enforces the \emph{one-sided} deficit mechanism: only a positive shortfall above the harvestable/share-adjusted flow liquidates stock or accrues debt; a surplus ($E \leq{} \sigma{}bA$, $E \leq{} B$) contributes nothing.
\item \textbf{$b_{G}$ is a conversion factor between the two "books", not a yield.} It is the value (in global hectares) of one hectare of standing stock, with units $gha\cdot{}ha^{-1}$. In (2) the capital-growth term $b_{G} G(A)$ has units $gha\cdot{}ha^{-1} \cdot{} ha\cdot{}yr^{-1} = gha\cdot{}yr^{-1}$, matching the flow term $bA$ ($b$ in $gha\cdot{}ha^{-1}\cdot{}yr^{-1}$ \ensuremath{\times{}} $A$ in $ha$), so $B$ is a flow and the two terms are commensurable. In (4) dividing the overshoot $[E - \sigma{}bA]_{+}$ (a $gha\cdot{}yr^{-1}$ flow) by $b_{G}$ ($gha\cdot{}ha^{-1}$) yields a \textbf{stock-loss rate} in $ha\cdot{}yr^{-1}$, so \texttt{1/b\_G} is the harvest coefficient $\gamma{} = 1/b_{G} = 1/V$ of §4.1 ($V$ = standing biomass \ensuremath{\div{}} annual production, \ensuremath{\approx{}} 20–100 yr). $b_{G}$ is therefore \textbf{not} dimensionless; it is the book-conversion rate that turns a flow deficit into capital loss.
\item \textbf{What debt degrades.} Equation (8) degrades \textbf{only the flow-yield coefficient $b$} (via $e^{-\alpha{}D}$): cumulative ecological debt erodes the surviving stock's \emph{productivity per unit area}, not the stock itself. It does \textbf{not} degrade $b_{G}$ (the standing-stock value), $\rho{}$ (the regeneration rate), or $A_{max}$ (the carrying-area ceiling); those remain fixed at their baseline values. This is Assumption (8) and is a deliberate scope: the model keeps the degradation channel (flow productivity) separate from the regeneration channel (rate and ceiling), so that the two do not double-count. A debt dependence on $G(A)$ or $b_{G}$ would be an additional, separately-argued extension.
\end{itemize}

\textbf{Delay structure.} The two delays enter \textbf{only} where stated. Regeneration is delayed, $G(A(t-\tau{}_{g}))$ in (4): the environment responds to its own state $\tau{}_{g}$ years earlier. Population response is delayed, $K(t-\tau{}_{p})$ in (5): the population adjusts to the carrying capacity it "knew" $\tau{}_{p}$ years ago. Debt formation (7) and harvest (via $E = eP$ in (4)) act on \textbf{current} values — the lags are in recruitment and demography, not in the depletion or debt terms. Crucially, $K = B/e$ (6) is \textbf{algebraic and not itself delayed}: carrying capacity is computed from the \emph{current} state $B = bA + b_{G} G(A)$, and the delay enters only through the population's response to a delayed value of it, $K(t-\tau{}_{p})$. The system is genuinely three-dimensional in the states $A, P, D$; $K$ carries no dynamics of its own. Equation (4) is the central structural assumption. Demand is met first from the flow yield $bA$ (the fruit — leave the tree); only the shortfall $[E - bA]_{+}$ liquidates the standing stock, at rate \texttt{1/b\_G}. This is the deficit mechanism, and it is why biocapacity appears in the stock equation: the two flow accounts are no longer parallel and unlinked. In the pure-capital limit ($\psi{} \to{} 0$, equivalently the saw-log form) equation (4) reduces to the standard gross-depletion harvest equation with harvest coefficient $\gamma{} = 1/b_{G}$. The model is well posed at the boundary. The original \texttt{dP/dt} diverges ($P/K \to{} \infty{}$ as $K \to{} 0$); we impose an extinction floor $A_{ext} > 0$ so that $K \geq{} b\cdot{}A_{ext}/e > 0$ is bounded and clamp $A \geq{} A_{ext}$, $P \geq{} 0$. "Collapse" is then a well-defined approach to the extinction boundary, not a numerically clamped blow-up (see §Numerics). The three state variables $A, P, D$ live in three distinct "books", and the model is written so these are never mixed: $A$ is the capital book (the standing base, removed only by liquidation); the flow yield $bA$ is the flow-yield book (the fruit, separable from the base); and $D$ is the degradation/debt book (the erosion of the surviving stock's yield). Conservation is a property of the incidence structure (each outflow is assigned to exactly one donor book) and positivity follows from donor limitation (an outflow vanishes when its donor compartment is empty). Any parameter pinned at an optimisation bound is reported as a declared fit defect rather than silently as a calibration. ---

\section{Assumptions}

Each equation carries an explicit modelling choice.

\begin{itemize}
\item \textbf{(1) Logistic regeneration} — symmetric self-limiting growth; a modelling convenience, not a law. The point \texttt{A\_max/2} is the maximal-growth point, not a threshold.
\item \textbf{(2) Biocapacity is additive flow + capital growth} — the flow is separable (crops, orchard); the capital growth is removable only by drawing down the stock (forest, fish). The composition is captured by the flow share $\psi{} = bA*/B*$, which is the master parameter separating the two ratios $R_{A}$ and $R_{B}$ (§4.5) and locating the masking illusion (prediction 5).
\item \textbf{(3) Per-capita footprint is constant} — $e$ (and the per-capita requirement it represents) is held constant as a baseline modelling choice, since this is the only way to isolate the stock–flow–demand feedbacks; endogenising $e$ as a time-varying per-capita requirement is an offered extension. The model does not carry a co-evolving per-capita requirement, and this is stated explicitly rather than left implicit.
\item \textbf{(4) Deficit-driven, immediate depletion} — liquidation is immediate; the delayed response is in recruitment $\tau{}_{g}$ (a tree takes $\tau{}_{g} \approx{} 20–80 yr$ to bear fruit), not in the depletion term.
\item \textbf{(5) Demographic delay on carrying capacity} — the population responds to the delayed carrying capacity $K(t-\tau{}_{p})$ (the conditions a cohort "knew"), Hutchinson-style.
\item \textbf{(6) $K$ is algebraic, not a state} — a function of the state; the system is genuinely 3-D ($A, P, D$).
\item \textbf{(7) Debt accrues only in overshoot, repayed at $\eta{}$} — $\eta{}$ is primary (it decides whether an equilibrium exists). We set $\eta{} = 0.05 yr^{-1}$, a minimal environmental regeneration / degradation-removal rate (slow, \ensuremath{\approx{}}20-yr timescale); the model is insensitive to $\eta{}$ over a modest range provided $\eta{} > 0$.
\item \textbf{(8) Degradation erodes the surviving stock's yield} — a separately-evidenced channel (soil fertility, over-picking), not double-counting the trees removed by $A$.
\item \textbf{(9) Technology is bounded} — a logistic wave; saturation is why the masking is transient.
\end{itemize}

\textbf{Delay asymmetry.} Depletion/liquidation is immediate while functional degradation is a separate, slower state. We justify this asymmetry explicitly and note that a third lag $\tau{}_{D}$ could be added ($dD/dt = [E(t-\tau{}_{D}) - B]_{+} - \eta{}D$) to keep the delay structure fully consistent. We also declare the aggregation of the governance clock: $\tau{}_{p}$ lumps several institutional time objects (observation, assessment, review, decision/deployment, ecological response, memory) and is not a single measured lag. ---

\section{Analytic results}

This section develops the analytic core in two registers that read the same object differently. The dynamical view (§4.1–§4.4, §Numerics, §Discussion) asks how the system behaves; the observational view (§4.5, §Predictions, §Discussion) asks what a policy-maker can watch. The thesis that holds them together is stated in §Presentation and in the abstract: the balance point is $R_{B} = 1$; $R_{A} = 1$ is a leading but non-causal signal; and when the regeneration lag is too long neither ratio warns — the lag, not the ratio, is the controlling variable.

\subsection{Carrying capacity, MSY, and the emergent ceiling}

$K = B/e$ is emergent (it moves with $A$ and $b$). Because the capital growth $b_{G} G(A)$ is positive only for $0 < A < A_{max}$, the total sustainable-yield curve $B(A) = bA + b_{G} G(A)$ has an interior maximum — the maximum sustainable yield (MSY): via $dB/dA = b + b_{G} \rho{}(1 - 2A/A_{max}) = 0$,

\begingroup\small\begin{verbatim}
A*  = A_max (b + b_G ρ) / (2 b_G ρ)
B_max = A_max (b + b_G ρ)² / (4 b_G ρ)
\end{verbatim}\endgroup

(this derivation is set out in full in §S3.1). So the sustainable equilibrium is interior ($A* < A_{max}$), not $A_{max}$ — a steady harvest of the capital growth keeps the stock below its unmanaged maximum. This is the surplus-production / MSY structure of Schaefer (1954), the same symmetric intrinsic-growth curve that gives $MSY = rK/4$. Only in the flow-only limit ($b_{G} G \ll{} bA$, the orchard) does the stock tend to $A_{max}$ with a zero harvest drain. \textbf{The interior-MSY statement is regime-conditional.} $A* < A_{max}$ holds only in the capital-dominated regime, i.e. when $b_{G} \rho{} > b$ (equivalently $\psi{} \to{} 0$). When $b_{G} \rho{} < b$ (the flow-dominated regime, $\psi{} \to{} 1$) $B(A)$ is monotone increasing on $[0, A_{max}]$ and there is no interior maximum: the sustainable point is the boundary $A_{max}$. This is precisely the regime the baseline parameter set sits in ($b_{G} \rho{} = 0.8\cdot{}0.05 = 0.04 < b = 0.5$), so the baseline sustainable state is the boundary $A_{max}$. The regime must be stated before quoting $A*$/$B_{max}$; they are the capital-only limit, not universal. \textbf{The flow share $\psi{} = bA*/B*$ interpolates between the two limit cases.} Report its regime dependence explicitly, as it governs which analytic core applies:

\begin{itemize}
\item $\psi{} \to{} 0$ (capital-dominated): interior $A*$; the stability classification of §4.3 applies.
\item $\psi{} \to{} 1$ (flow-dominated, the orchard limit): $A_{max}$ at sustainability, a liquidation mask, and the vicious cycle — the $a_{11} = \rho{}(1-2A*/A_{max}) + b/b_{G}$ liquidation feedback.
\item $0 < \psi{} < 1$: interior $A*$ and liquidation feedback coexist. This is the general object, and the increasing-harvest and orchard models are recovered as limits.
\item $\gamma{} = 1/b_{G} = 1/V$, the orchard's salvage value ($V$ = standing biomass \ensuremath{\div{}} annual production, \ensuremath{\approx{}} 20–100 yr). The capital-harvest coefficient is therefore measurable and is not a free parameter.
\item $\psi{}$ is a testable prediction (prediction 5). Small $\psi{}$ (flow yield is a small share, cropland/grazing) puts the model near the capital limit, where a smooth decline is stabilised by demographic feedback and the illusion is small. Large $\psi{}$ (forest/fishery) makes the overshoot invisible until demand exceeds the flow yield, then liquidation with a threshold $A_{c}(E)$. The model therefore predicts the masking illusion is more visible in flow-dominated (high-$\psi{}$) systems than in capital-dominated (low-$\psi{}$) ones.
\end{itemize}

\textbf{Regime \ensuremath{\to{}} outcome.} The single testable condition $b_{G} \rho{} \lesseqgtr{} b$ (equivalently $\psi{} \lesseqgtr{} 1/2$) decides everything downstream:

\begin{table}[htbp]
\centering
\small
\begin{tabularx}{\textwidth}{>{\raggedright\arraybackslash}X>{\raggedright\arraybackslash}X>{\raggedright\arraybackslash}X>{\raggedright\arraybackslash}X>{\raggedright\arraybackslash}X>{\raggedright\arraybackslash}X}
\toprule
Regime & Condition & $B(A)$ shape & Sustainable point & MSY & Collapse/mask behaviour \\
\midrule
capital-dominated & $b_{G} \rho{} > b$ ($\psi{} \to{} 0$) & interior maximum & interior $A* < A_{max}$ & interior MSY $A*$, $B_{max}$ & smooth decline stabilised by demographic feedback — illusion small \\
flow-dominated (baseline) & $b_{G} \rho{} < b$ ($\psi{} \to{} 1$) & monotone \ensuremath{\uparrow{}} & boundary $A_{max}$ & none interior (max at $A_{max}$) & liquidation with threshold $A_{c}(E)$; overshoot invisible until demand exceeds flow yield — illusion more visible \\
marginal & $b_{G} \rho{} = b$ & max at the boundary & $A* = A_{max}$ (interior max merges with boundary) & boundary $A_{max}$ & $\psi{} \to{} 1$ \\
\bottomrule
\end{tabularx}
\end{table}

A parameter sitting exactly on the $b_{G} \rho{} = b$ boundary is a non-generic choice and should be stated and perturbed off.

\subsection{Fixed-liability threshold and the saddle-node = MSY}

For a fixed liability $E$, the deficit-region equation $dA/dt = 0$ is a quadratic in $A$ with an unstable root (the separatrix), and the saddle-node is exactly the MSY:

\begingroup\small\begin{verbatim}
A_c(E) = [ (b_G ρ + b) − √((b_G ρ + b)² − 4 b_G ρ E/A_max) ] A_max / (2 b_G ρ)
E_sn   = B_max = A_max (b + b_G ρ)² / (4 b_G ρ)            (safe basin vanishes at MSY)
\end{verbatim}\endgroup

(the fixed-liability threshold is derived in §S3.2). A fixed liability above $B_{max}$ cannot be sustained by harvesting regrowth. The threshold is emergent and $E$-dependent (no Allee term needed), located at $A_{c}(E)$, never at \texttt{A\_max/2}. When realised, the fold is a saddle-node (fold) bifurcation — the safe basin vanishes as $E$ crosses $B_{max}$, the catastrophic-shift / critical-transition structure of alternative-stable-state ecology (Scheffer, Carpenter, Foley, Folke \& Walker 2001): a loss of resilience precedes the switch, and the switch is abrupt and hard to reverse, not because of an imposed Allee minimum but because the stable and unstable equilibria first meet at the fold and then destroy each other. \textbf{The fold is regime-scoped.} It is realised only in the capital-dominated regime ($b_{G} \rho{} > b$), where $B(A)$ has an interior maximum. There $dB/dA = 0$ exactly at the interior MSY, and for $E$ just below $B_{max}$ there are two fixed points — one unstable ($B'(A) > 0$) and one stable ($B'(A) < 0$) — that merge at the saddle-node (verified numerically: $B'(A*) = 0$; two roots with opposite-sign $B'$ for $E < B_{max}$). At baseline ($b_{G} \rho{} = 0.04 < b = 0.5$), $B(A)$ is monotone increasing on $[0, A_{max}]$ with $B'(A) > 0$ everywhere, so no interior fold is realised: the fixed-liability equilibria are single-valued and unstable, and the stable object is the boundary $A_{max}$. \textbf{The fold must not be conflated with the $\tau{}_{g}$ cliff.} As we show in §Numerics and §Discussion, the $\tau{}_{g}$ transition is a basin-boundary crisis, not a fold: the leading eigenvalue is constant and positive ($Re \lambda{} \approx{} +0.62$) for every $\tau{}_{g}$, so no eigenvalue crosses zero as $\tau{}_{g} \to{} 20$ and there is no critical-slowing-down / early-warning precursor (the $P$-relaxation return time is flat \ensuremath{\approx{}}150 yr for $\tau{}_{g} = 0–18$, then the attractor vanishes abruptly rather than slowing). Scheffer's "loss of resilience precedes the switch" therefore describes only a capital-dominated $E$-fold; it is the wrong description of the delay transition. The two must be told apart (§Discussion).

\subsection{Stability; the vicious cycle; the stability classification}

\textbf{The deficit = stock-decline identity.} When demand exceeds the flow yield ($E > bA$), Eq. (4) reduces exactly:

\begingroup\small\begin{verbatim}
dA/dt = G(A(t−τ_g)) − [E − bA]₊/b_G  =  (B̃ − E)/b_G
\end{verbatim}\endgroup

because $[E - bA]_{+} = E - bA$ there, and the lag-adjusted biocapacity is $\tilde{B} = bA + b_{G} G(A(t-\tau{}_{g}))$ (so $\tilde{B} = B$ only when regeneration is evaluated at the delayed state). This is the algebraic identity that makes "deficit-driven" rigorous: the shortfall between biocapacity and footprint, not the gross harvest, drives stock decline, and the model's stock build-up exactly cancels against its flow term. Corollary (in the $R_{B}$ form of §4.5): $dA/dt < 0 \iff{} E > \tilde{B}$, and at equilibrium $E = \tilde{B}$, so $R_{B} = E/\tilde{B} = 1$ at the balance point. The $R_{B} = 1$ locus of §4.5 is exactly the neutral continuum on which the characteristic equation has $D(0) = 0$ (§Numerics) — the separator and the continuum are one and the same object, seen from the ratio (monitoring) and spectral (dynamical) views. Linearising (4)+(5) in the deficit region gives (at the interior point, under \ensuremath{\sigma{}} = 1 so that the equilibrium family $P = B(A)/e$ is well defined)

\begingroup\small\begin{verbatim}
a₁₁ = ρ(1 − 2A*/A_max) + b/b_G      ;  a₁₂ = −e/b_G      ;  a₂₂ = −r
det = r·ρ·A*/A_max > 0  (never a saddle)   →   zero-delay condition is  a₁₁ < r  (NOT a₁₁ < 0)
\end{verbatim}\endgroup

\textbf{Sufficient condition for the monotone instability (the "positive real eigenvalue for every delay" claim).} Linearising gives the two-delay characteristic equation $D(s;\tau{}_{g},\tau{}_{p}) = (s - a_{1}e^{-s\tau{}_{g}} - a_{3})(s + r) - a_{E} a_{4} e^{-s\tau{}_{p}} = 0$ with $a_{1} = G\prime{}(A*)$, $a_{3} = b/b_{G}$, $a_{E} = -e/b_{G}$, $a_{4} = r K\prime{}(A*)$, and the equilibrium condition $P* = K(A*)$ forces $-a_{E} a_{4} = r\cdot{}S$, where $S = a_{1} + a_{3} = G\prime{}(A*) + b/b_{G}$. Then:

\begin{itemize}
\item $D(0) = 0$ on the whole family — the neutral continuum, so there is \textbf{no isolated interior attractor}.
\item With no delay the non-neutral root is $\lambda{} = S - r$. Hence a \textbf{positive real root} exists at zero delay exactly when $S > r$, i.e. $G\prime{}(A*) + b/b_{G} > r$.
\item The root persists \textbf{for every delay} provided $F\prime{}(0) = r - S + r a_{1} \tau{}_{g} - r S \tau{}_{p} < 0$. At baseline ($a_{1} = -0.0167$, $S = 0.6083$, $r = 0.02$) this holds over the whole $(\tau{}_{g}, \tau{}_{p}) \in{} [0,400]²$ plane (max $F\prime{}(0) = -0.5883$), so a positive real eigenvalue is guaranteed for every delay combination (verified: leading $Re \lambda{} = +0.625$ real across the plane; the exact $s = i\omega{}$ crossing-curve scan finds zero imaginary-axis crossings).
\end{itemize}

\textbf{So the claim is conditional, not parameter-free.} A positive real eigenvalue for every delay is guaranteed only when $S = b/b_{G} + G\prime{}(A*) > r$ (and the delay condition above); it is \emph{not} a universal theorem for arbitrary parameter values. It is stated here for the representative parameter set, as checked numerically in §Numerics and §S5. Under a reservation $\sigma{} < 1$ the equilibrium locus is no longer $P = B(A)/e$ but $P = (\sigma{}bA + b_{G} G(A))/e < B/e$, and the whole family (and hence the neutral continuum) is shifted — the $\sigma{} = 1$ restriction is required for the results of §4.3 and §4.5 as stated (verified numerically: at $\sigma{} = 0.9$, $dA/dt(P = B/e)$ is negative for every interior $A*$).

\begin{itemize}
\item \textbf{The vicious cycle is real and quantitative.} $a_{11}$ gains the \texttt{+b/b\_G} term. For liabilities that scale with the stock ($E = f\cdot{}bA$) the interior point is $A* = A_{max}[1 - (f-1)b/(b_{G} \rho{})]$, and it is self-sustaining (the environment's own mode goes locally runaway, $a_{11} > 0$) precisely when $(2f-1)\nu{} > 1$, where $\nu{} = b/(b_{G} \rho{})$. With $b_{G}$ the standing-stock value (20–100 yr) and realistic $\rho{} \approx{} 0.02–0.1 yr^{-1}$, $\nu{}$ is $O(0.1–2)$; treat $(2f-1)\nu{} > 1$ only as an estimated indicator, not a measured threshold, because $f$ and the liability elasticity are set by the model. The "vicious cycle" therefore describes Eq. (4) and not a gross-depletion form.
\item \textbf{Stability classification (sign-corrected).} For $\rho{} \gg{} r$, the 2-D two-delay system reduces to a scalar two-gain delayed logistic with control $\chi{} = q/(\rho{} - 2q)$, $q = \gamma{} e b_{0}/r_{opt}$ (depletion pressure): $\chi{} > 1$ ($\Lambda{} > 0$) \ensuremath{\Rightarrow{}} $\tau{}_{g}$-only Hopf: $\omega{} = r\sqrt{}(\chi{}²-1)$, $\tau{}_{g}* = arccos(-1/\chi{})/\omega{}$. $\chi{} < 1$ ($\Lambda{} < 0$) \ensuremath{\Rightarrow{}} $\tau{}_{p}$-only Hopf: $\omega{} = r\sqrt{}(1-\chi{}²)$, $\tau{}_{p}* = arccos(-\chi{})/\omega{}$. $\chi{} = 1$ ($\Lambda{} = 0$, i.e. $\rho{} = 3q$) \ensuremath{\Rightarrow{}} neither; two-delay boundary $s = \pi{}/\omega{}$, $\omega{} = 2r\cdot{}cos(\omega{}d/2)$ \ensuremath{\Rightarrow{}} $s \approx{} \pi{}/(2r) = 78.5 yr$ at $d \to{} 0$. The sign structure is exact (from the $|Q| = |P|$ elimination); only $\omega{}*, \tau{}*$ are $O(r/\rho{})$ approximate, and the fast–slow reduction is valid only while $a_{11} \leq{} 0$ (else use the full 2-D transcendental equation, where the $A$-mode can itself oscillate). On the constant-parameter S0 of the present model this Hopf classification is not realised: the exact $[\cdot{}]_{+}$ gives a monotone positive-real eigenvalue with no imaginary-axis crossing for every delay (§Numerics). The classification above describes the fast–slow reduction and does not predict the actual (structural) instability. \textbf{Baseline sits at a knife-edge.} Baseline $\chi{} = 1$ because $\rho{} = 3q$, i.e. $\Lambda{} = 0$, the measure-zero surface where the two gain surfaces $r²a_{11}² = (\gamma{} e a_{21})²$ are equal ($g_{M} = g_{P}$). This is a non-generic, unexplained parameter choice, and it is precisely why "no single delay destabilises" holds at baseline. We do not justify it; we report instead the generic result that follows from the sign of $\Lambda{}$: perturbing off the knife-edge gives exactly one lag destabilising — $\tau{}_{g}$-only if $\chi{} > 1$ ($\Lambda{} > 0$), $\tau{}_{p}$-only if $\chi{} < 1$ ($\Lambda{} < 0$). Any scan range and any "no single delay destabilises" claim must be stated as a function of $\Lambda{}$/$\chi{}$, never as a parameter-free statement (§Numerics).
\end{itemize}

\subsection{The complete dimensionless group set}

$\chi{} = q/(\rho{}-2q)$ is the cleanest single control (it matches the scalar two-gain picture), but the paper carries the complete non-dimensionalization as the full generality statement: $\hat{t} = rt$, $a = A/A_{max}$, $p = P\cdot{}r_{opt}/(b_{0}A_{max})$, and the groups $s = \rho{}/r$, $g = \gamma{} b_{0} f/\rho{}$, $f = e/r_{opt}$, $\theta{}$, plus scaled delays $\hat{\tau{}}_{M} = r\tau{}_{g}$, $\hat{\tau{}}_{P} = r\tau{}_{p}$ (full detail and scaling in §S2). We present both: $\chi{}$ for the clean stability-sign rule, and the six-group set for full generality — they are complementary, not competing.

\subsection{Macro-ratio safe operating space ($R_{B}$ vs $R_{A}$)}

In addition to the dimensionless groups, the bookkeeping can be reported as two observable ratios that summarise the safe-operating picture. Define the biocapacity ratio and the flow-yield ratio:

\begingroup\small\begin{verbatim}
R_B = E / B            (footprint ÷ total biocapacity)
R_A = E / (b A)        (footprint ÷ flow yield)
ψ   = bA / B           (flow share, 0 ≤ ψ ≤ 1)
R_A = R_B / ψ
\end{verbatim}\endgroup

\textbf{The operating boundary is $R_{B} = 1$, and it is exactly the neutral equilibrium family of §4.3 $P = B(A)/e$.} From the deficit identity $dA/dt = (\tilde{B} - E)/b_{G}$ (§4.3), $dA/dt < 0 \iff{} E > \tilde{B}$, so the stock declines iff footprint exceeds total biocapacity. At equilibrium $E = \tilde{B}$ so $R_{B} = 1$; this is simultaneously the bookkeeping trigger, the macro-ratio safe-operating boundary, and (as the family $P = B(A)/e$) the basin separator. Verified onset (baseline S0, documented $A_{0}$ grid): the recover/collapse boundary is crossed at $R_{B} = 1.000$ for every $A_{0} \in{} {0.30,\ldots{},1.05}$, while the flow-yield ratio there is $R_{A} = 1.01–1.06$. \textbf{$R_{A} = 1$ is necessary but not sufficient.} $R_{A} = 1$ means $E = bA$, i.e. demand equals the flow only. A system sitting between $R_{A} = 1$ and $R_{B} = 1$ still regenerates ($dA/dt > 0$), because the capital growth $b_{G} G(A)$ is untouched. $R_{A} > 1$ is therefore not the collapse trigger, and the two ratios are not interchangeable: $R_{A}$ is a leading, non-causal signal; $R_{B}$ is the trigger. \textbf{The gap between them is the flow share, which is the master parameter.} $R_{A} = R_{B}/\psi{}$, so the two ratios coincide in the flow-dominated limit ($\psi{} \to{} 1$, baseline) and separate as $\psi{} \to{} 0$ (capital-dominated). Regime-scoped closed form (at a sustainable interior MSY): $\psi{}* = 2/(1 + b_{G}\rho{}/b)$, $R_{A}^eq = (1 + b_{G}\rho{}/b)/2$. Verified across regimes: $b_{G}\rho{}/b = 0.05$ (flow-dominated, no interior MSY, $\psi{} \to{} 1$, $R_{A} = 1$); $2.0 \to{} \psi{}* = 0.667, R_{A}^eq = 1.50$; $37.5 \to{} \psi{}* = 0.052, R_{A}^eq = 19.25$; $150 \to{} \psi{}* = 0.013, R_{A}^eq = 75.5$. The single regime index $b_{G}\rho{}/b$ therefore governs the gap between $R_{A}$ and $R_{B}$ ($\psi{}* = 1/R_{A}^eq = 2/(1 + b_{G}\rho{}/b)$). \textbf{A sustainable system can run at $R_{A} > 1$.} In the capital-dominated regime the balance point is $R_{B} = 1$ with $R_{A} > 1$ (footprint > flow yield, the buffer intact). $R_{A} > 1$ is not a warning; the balance point is $R_{B} = 1$ in every regime. \textbf{Temporal-lead nuance.} Because $R_{A} = R_{B}/\psi{}$, $R_{A}$ crosses 1 earlier than $R_{B}$ by the factor $1/\psi{}$, but its lead is \ensuremath{\approx{}}0 in flow-dominated systems and grows only with capital dominance. A large lead is therefore a signature of the capital-dominated regime, not a general warning property. \textbf{Fold \ensuremath{\leftrightarrow{}} $R_{B} = 1$ (regime-scoped).} The §4.2 fixed-liability fold threshold $E_{sn} = B_{max} = A_{max}(b+b_{G}\rho{})²/(4b_{G}\rho{})$ is precisely the $R_{B} = 1$ boundary evaluated at $(A*, E = B_{max})$. Verified (capital-dominated $b = 0.02$): interior MSY $A* = A_{max}(b+b_{G}\rho{})/(2b_{G}\rho{}) = 0.900$, $B(A*) = B_{max} = 0.0270$, $E_{sn} = 0.0270$, so $R_{B} = E_{sn}/B_{max} = 1.000$ exactly — the fold and the balance point are the same object. This fold is realised only when $b_{G}\rho{} > b$, the identical scoping as the §4.1 regime table and §4.2; at baseline it is absent. \textbf{The separator degrades with the lag — two sides of one coin.} At short/no delay the neutral family $P = B(A)/e$ (i.e. $R_{B} = 1$) separates recover from collapse at balanced accuracy 99.2 \% (no-delay) \ensuremath{\to{}} 81.2 \% ($\tau{}_{g} = 30$); a linear functional is chance at $\tau{}_{g} = 30$ (balanced 50.0 \%, raw 94.7 \% = the majority-collapse class, a class-imbalance artefact). So $R_{B} = 1$ is necessary for recovery but not sufficient once the lag is too long — the operating boundary (here) and the basin crisis (§Numerics, §Discussion) are complementary statements about one object. The family $P = B(A)/e$ is the separator, not a "first integral". \textbf{One object, many faces.} The single scalar $R_{B} = 1$ is simultaneously the bookkeeping trigger ($dA/dt < 0 \iff{} E > \tilde{B}$, §2.2, §4.3); the basin separator $P = B(A)/e$ (§4.3, §Numerics — the neutral continuum); the interior-MSY fold threshold of §4.2, only when $b_{G}\rho{} > b$ (absent at baseline); the target of prediction 8 (§Predictions); and the object whose long-$\tau{}_{g}$ failure produces the silent-collapse and measure-zero-rescue negative results of §Discussion. No new parameter appears; each face is the same statement "$E = B$" in a different coordinate. \textbf{Figures (observational view).} \texttt{scans/topdown\_macro\_ratios.png} — the macro-ratio plane ($R_{B}$ vs $R_{A}$, two panels $\tau{}_{g}=10$ and 30), showing the $R_{B}=1$ boundary as both the basin separator (recover/collapse) and the safe-operating boundary, with the $R_{A}>1$ buffer. \texttt{scans/topdown\_ratio\_separation.png} — the \ensuremath{\psi{}}-regime closed form $R_{A}^eq=(1+b_{G}\rho{}/b)/2$ (and $\psi{}*=2/(1+b_{G}\rho{}/b)$), i.e. how far the leading indicator $R_{A}$ sits above the trigger $R_{B}$ as the regime index $b_{G}\rho{}/b$ grows. (The $\tau{}_{g}$ cliff and full-plane no-Hopf belong to the dynamical view and are in §Numerics.) ---

\section{Results: principal claims}

The results established by the analysis are the following; each is grounded in §4, §Numerics, and §Discussion.

\begin{itemize}
\item \textbf{The constant-parameter subsystem is monotonically unstable, not Hopf-driven.} Under full harvest ($\sigma{} = 1$), so that the equilibrium family $P = B(A)/e$ is well defined, and for the representative parameter set it has a positive real eigenvalue (\ensuremath{\approx{}} +0.62) for every delay, with no imaginary-axis crossing — guaranteed by the sufficient condition $S = b/b_{G} + G\prime{}(A*) > r$ (here $0.6083 > 0.02$), not as a parameter-free theorem; and $D(0)=0$ on the whole equilibrium family, so there is no isolated interior attractor to frame a delay-ratio Hopf. This contrasts with the classical delayed-logistic Hopf of Hutchinson-type single- and two-delay models, in which a delay ratio produces an oscillatory onset; here the onset is a structural vicious cycle. (The $78.5 yr = \pi{}/(2r)$ coincidence is a two-loop coincidence, not the Hutchinson $\pi{}/(2r)$ demographic threshold.)
\item \textbf{The productivity illusion is real but conditional and narrow.} A genuine $B$-rising-while-$A$-falls window exists only for a small initial deficit (\ensuremath{\approx{}}5.4 yr at $E-b_{0}A_{0}=0.06$, collapsing to zero at deficit \ensuremath{\approx{}}0.075; converged RK4 — see §Demonstration). It is narrow, deficit-bounded, and transient, with a computable $t_{peak}$ and a critical deficit.
\item \textbf{Technology can increase cumulative debt.} Under overshoot, yield technology raises $K$, $P$ and $E$, and can increase cumulative debt — a Jevons-type rebound.
\item \textbf{The true threshold is emergent.} There is no threshold at \texttt{A\_max/2}. The real threshold is emergent: $A_{c}(E)$ (fixed liability) or the MSY $B_{max}$; \texttt{A\_max/2} is only the maximal-regeneration point.
\item \textbf{The full overshoot model has an equilibrium only because of $\eta{}$.}
\item \textbf{The full system is three delayed states ($A,P,D$); only the $\alpha{}=0$ subsystem is 2-D; $K$ is algebraic.}
\item \textbf{The operating (decline) trigger is $R_{B} = 1$ (footprint = total biocapacity), not $R_{A} = 1$.} $dA/dt < 0 \iff{} E > \tilde{B}$; $E > bA$ ($R_{A} = 1$) is necessary-but-not-sufficient. A system can sit on the collapse side with $R_{A} > 1$ (capital-dominated) or, at long $\tau{}_{g}$, with both ratios below 1 (silent collapse, §Discussion).
\item \textbf{The debt-compounding asymmetry is a theorem under multiplicative degradation.} Under $b = (b_{0}+T_{b})e^{-\alpha{}D}$, "debt compounds without bound while technology saturates" follows from the model; it is false under an additive form.
\end{itemize}

\textbf{Scenario B/C is the inverse of the orchard framing.} In Scenarios B and C the environmental stock $A$ rebounds to \ensuremath{\approx{}}$A_{max}$ (\ensuremath{\approx{}}1.19) while population $P$ and the harvest/biocapacity $B$ collapse — an "environment recovers, humans collapse" outcome that is the opposite of the orchard framing ("humans die, orchard survives"). This is a scenario outcome, not a claim to be over-claimed. ---

\section{Falsifiable predictions}

These are the testable content of the model. They fall into two registers: predictions 1–3, 6 and 7 are dynamical (stability, onset, recovery dynamics); prediction 5 straddles the two views (the \ensuremath{\psi{}}-mediated mask, §4.5); and prediction 8 is the observational one (the ratio safe-operating space and whether a monitor can forecast collapse, §4.5). The mapping of each prediction to the section that establishes it is tabulated in §S6.

\begin{enumerate}
\item \textbf{Which lag destabilises} is set by the sign of $\Lambda{}$/$\chi{}$: on the constant-parameter subsystem (under full harvest, $\sigma{} = 1$) the interior point is monotonically unstable for every delay — provided $b/b_{G} + G\prime{}(A*) > r$, the representative-case condition — giving a positive real eigenvalue and no imaginary-axis crossing, so the "which-lag" question has no Hopf answer there.
\item The \textbf{oscillation period near onset \ensuremath{\approx{}} 4\ensuremath{\times{}} the dominant lag} — applies only on a Hopf boundary; the constant-parameter subsystem has no oscillatory onset (monotone collapse), though the full model (with $\eta{}$, $D$) can still show damped oscillatory transients.
\item The productivity illusion has a \textbf{computable peak-biocapacity time $t_{peak}$} with the signature "$B$ rising while $A$ falls".
\item \textbf{Reducing a policy lag $\tau{}_{e}$ matters comparably to reducing the overshoot $f$} — a testable, policy-relevant ranking.
\item \textbf{The flow-share $\psi{}$ locates the masking illusion} (§4.1, §4.5): the "$B$ rises while $A$ falls" window should be more visible in flow-dominated (high-$\psi{}$, $b_{G}\rho{} < b$) systems than in capital-dominated (low-$\psi{}$, $b_{G}\rho{} > b$) ones.
\item \textbf{The regeneration lag $\tau{}_{g}$ sets recovery vs collapse} (§Numerics): the recover fraction collapses from \ensuremath{\approx{}}40 \% to \ensuremath{\approx{}}5 \% across a steep \textbf{18–20 yr band}, so the model predicts collapse whenever the regeneration lag exceeds \ensuremath{\approx{}}20 yr and recovery only for $\tau{}_{g} \lesssim{} 18$. The empirically anchored effective timescale from the published recovery statistics is \ensuremath{\approx{}}25–33 yr (representative \ensuremath{\approx{}}29 yr; SI §S5.1) — at or above the collapse threshold — with primary forests at the top of that range and many soils and (non-recovered) fisheries in or beyond it. This is the necessary-but-not-sufficient reading of §4.5: $R_{B} = 1$ is necessary but not sufficient once $\tau{}_{g}$ is long. It is a quantitative, direct prediction tied to the empirical $\tau{}_{g}$ band.
\item \textbf{Regeneration \emph{rate} vs \emph{lag} are two decoupled levers} (§Discussion): time-to-50 \% scales \ensuremath{\approx{}} inversely with the regeneration rate $\rho{}$, but whether the stock recovers at all is set by the regeneration lag $\tau{}_{g}$. A system with a long $\tau{}_{g}$ cannot be rescued by raising $\rho{}$ alone.
\item \textbf{The operating boundary is $R_{B} = 1$ (footprint = total biocapacity), not $R_{A} = 1$ (§4.5).} $R_{A}$ crosses 1 earlier than $R_{B}$ (by $1/\psi{}$) but is not the collapse trigger, and its lead is \ensuremath{\approx{}}0 in flow-dominated systems. A system can be on the collapse side while $R_{A} > 1$ (capital-dominated) — and, in the long-$\tau{}_{g}$ regime, even while $R_{B} < 1$ (silent collapse, §Discussion). Falsifiable: monitor the two ratios and confirm the trigger is $R_{B} = 1$, not the earlier-crossing $R_{A}$.
\end{enumerate}

\textbf{The two masks.} \emph{Technology mask:} $B$ rises while $A$ falls (requires progress). \emph{Liquidation mask:} $E$ steady while $A$ falls (requires none). The discriminating observable is exactly which holds. \textbf{Why the illusion is necessarily transient.} From $B = bA$ (flow component) the growth rate of biocapacity is $d ln B/dt = d ln b/dt + d ln A/dt$. Under the mask window $A$ falls, so the second term is negative; the first term is bounded because technology saturates ($T_{b}$ is a logistic wave, so $d ln b/dt \leq{} 0$ eventually). The sign of $d ln B/dt$ is therefore pinned by the unboundedly negative $d ln A/dt$ under the stock-liquidation cycle, so a $B$-rising-while-$A$-falls window cannot persist: the yield gain from a bounded technology wave is finite, whereas the stock decline it is meant to offset is unbounded once liquidation takes over. This is the analytic reason the mask is narrow and deficit-bounded (§Demonstration), and why the "productivity illusion reads as if sustainability were improving" is a transient, recover-to-reckoning signal rather than a new equilibrium. ---

\section{Presentation}

\begin{itemize}
\item \textbf{Thesis (stated once here and in §4):} the balance point is $R_{B} = 1$ (footprint = total biocapacity); the flow-yield ratio $R_{A} = 1$ is a leading but non-causal signal; and when the regeneration lag is too long neither ratio warns — the lag, not the ratio, is the controlling variable.
\item \textbf{Feedback diagram with a switch} (fruit harvest vs. capital liquidation): $A \to{}^{b} B \to{} K \to{} P \to{} E$, with $E - bA$ deciding the switch, $D \to{}^{\alpha{}} b$, and an exogenous bounded $T_{b}$. Two positive-feedback loops compound: Loop 1 (stock-liquidation: $A\downarrow{} \to{} B\downarrow{} \to{} deficit\uparrow{} \to{} liquidation\uparrow{} \to{} A\downarrow{}$) and Loop 2 (debt-erosion: $D\uparrow{} \to{} b\downarrow{} \to{} B\downarrow{} \to{} deficit\uparrow{} \to{} D\uparrow{}$). Rendered: \texttt{scans/feedback\_diagram.png}.
\item \textbf{Mappings.} In the orchard framing, the flowing fruit corresponds to biocapacity, and the standing trees correspond to the capital stock that generates it; the model's flow share $\psi{}$ is the share of sustainable biocapacity that comes from the flow.
\item \textbf{Non-smoothness and one-sided stability.} At the switch $E = bA$ the exact $[\cdot{}]_{+}$ is non-smooth, so the sustainable point is a boundary of the deficit regime, not an interior point. Report one-sided stability (linearise from the deficit side only). Because the exact switch yields no imaginary-axis crossing (§Numerics), the non-smoothness does not create spurious oscillations.
\item \textbf{Parameter justification.} The representative values are illustrative and must be caveated: $\rho{}$ is large in the original scale, and $\gamma{}$, $\alpha{}$, $\tau{}$ are lumped/effective parameters, not independently measured.
\item \textbf{Limitations of National Footprint Accounts.} The accounts are account-based (measured yields and conversion factors) and do not capture soil erosion, deforestation, or groundwater depletion, so estimated biocapacity likely overstates and the Footprint likely understates real overshoot (§Model scope).
\item \textbf{Units.} Symbols and units are given for $A_{max}$, $b_{0}$, $\Delta{}b$, and every other symbol in the symbol table (§2.1); no dimensionless parameter is left without its stated unit.
\end{itemize}

\textbf{Related literature.} The delayed-logistic \emph{stability} result is Hutchinson (1948), the classical single-delay logistic with its $\pi{}/(2r)$ demographic threshold. Haberl \& Aubauer (1992) \emph{apply} this delayed-logistic framework to human population/load dynamics; their contribution is the application, not the introduction of the time delay into population dynamics, which is Hutchinson. Brander \& Taylor (1998) is a distinct Ricardo–Malthus renewable-resource model. The surplus-production/MSY curve used here is Schaefer (1954), and the fold/saddle-node with a vanishing safe basin is the catastrophic-shift structure of Scheffer et al. (2001); the empirical regeneration-lag band follows Poorter et al. (2016), Poeplau et al. (2011), Hutchings \& Reynolds (2004), and Neubauer et al. (2013). On the data side, the National Footprint Accounts are documented in Wackernagel \& Rees (1996), Wackernagel et al. (2002), Borucke et al. (2013), and Lin et al. (2018). These accounts are not uncriticised: Blomqvist et al. (2013), Giampietro \& Saltelli (2014), and van den Bergh \& Grazi (2015) raise substantive methodological objections, and the Global Footprint Network-side rejoinders appear in Galli et al. (2016). These attribution statements are given so the reader is not misled about what each source establishes. ---

\section{Numerics, verification, and well-posedness}

\begin{itemize}
\item \textbf{Solver.} Method-of-steps with RK4 (or $dde23$/$pydelay$; Shampine \& Thompson 2001), the exact $[\cdot{}]_{+}$ switch as the model states it (never a smooth ramp, because a softplus ramp leaks depletion when $E < bA$), clamping $A \geq{} A_{ext} > 0$, $P \geq{} 0$, $K \geq{} K_{min}$, and an explicit extinction floor (so "collapse" is a model result, not a clamp). Where a smooth variant of the $[\cdot{}]_{+}$ switch is used (reference implementation only), fix its width $w$ and report insensitivity to it. Full-$(7)$ runs use $\eta{} = 0.05 yr^{-1}$; the constant- parameter subsystem S0 drops $D$ and $\eta{}$, so $\eta{}$ does not enter it.
\item \textbf{Protocol and convergence.} State $dt$, history functions, and the $\Delta{}t$-convergence table (e.g. scenario A $A*$: 0.8022 at $dt=0.5$ vs. 0.8021 at $dt=0.05$). Each reported number carries its protocol.
\item \textbf{Method sensitivity and numerical checks.} The final yield $b_{final} = 0.317$ gives $D \approx{} 6.76$. The endpoint $D_{E}$ is method-dependent (5.26 / 6.74 / 18.70 for the frozen/crashed/un-clamped endings), so it must be reported with its termination convention. The leading Hopf delays are $\tau{}_{g}* = 85.4 yr$ and $\tau{}_{p}* \approx{} 231 yr$.
\item \textbf{Measured basin shrinkage.} The qualitative "basin of attraction shrinks" is quantified. The stable fraction of the $(A_{0},P_{0})$ initial-condition plane for the overshoot subsystem falls from 0.506 (no delays) to 0.042 (baseline $(\tau{}_{g},\tau{}_{p})=(30,25)$), and the standard IC $(1.0,0.1)$ flips from stable to collapse. Report the stable fraction as a function of $(\tau{}_{g},\tau{}_{p})$; complement it with the closed-form separatrix criterion $A_{c}(E)$ (§4.2) — the measured fraction is the empirical witness, $A_{c}(E)$ the analytic boundary.
\item \textbf{Contrast with the gross-harvest formulation.} The deficit-driven model (Eq. 4) is structurally distinct from the standard gross-harvest (surplus-production) formulation in which total demand is removed directly from the stock. The gross-harvest subsystem has a unique interior attractor; the deficit-driven subsystem instead has a one-sided boundary at $A \to{} A_{max}, P \to{} b_{0}A_{max}/e$ and no robust interior attractor, so any population overshoot into $E > bA$ triggers the vicious-cycle collapse ($A \to{} A_{ext}$). This is the mechanistic reason the two formulations differ (see §Comparison of formulations).
\item \textbf{Basin recompute (constant-parameter subsystem).} On the documented $A_{0}\times{}P_{0}$ grid (13\ensuremath{\times{}}16 = 208; $P$ straddling $B* = b_{0}A_{max}/e \approx{} 1.09$), the recover fraction falls from 39.9 \% (no delay) to 5.3 \% (baseline $(30,25)$); the collapse fraction rises 60.1 \% \ensuremath{\to{}} 94.7 \%. The regeneration delay triggers an $A_{max}$-overshoot \ensuremath{\to{}} $E>bA$ \ensuremath{\to{}} liquidation collapse, and the recover basin vanishes abruptly as $\tau{}_{g}$ crosses \~{}20 yr. A time-domain illustration (\texttt{scans/r1\_recovery\_vs\_collapse.png}) shows the same IC $(A_{0},P_{0})=(1.0,0.1)$ recovering ($A\to{}A_{max}$, $\tau{}_{g}=10$) versus collapsing ($A\to{}A_{ext}$, overshoot $A\to{}1.36$, $\tau{}_{g}=30$); a recovery-side figure (\texttt{scans/eco\_recovery\_insight.png}) shows the density-dependent, sigmoidal recovery trajectory and the decoupling of regeneration rate (sets recovery time) from regeneration lag (sets recovery outcome).
\item \textbf{Scenario-D threshold is near-critical, not a clean result.} $(20,20)$ gives min $A = 0.631$ and recovers, while $(30,25)$ gives min $A < 0.6$ and collapses ($0.6 = A_{max}/2$, the maximal-regeneration point, not a threshold). State the actual basin boundary, not a single point.
\item \textbf{Reporting rigour.} (i) State the grid range of the stability scan and normalise "barely positive" $Re \lambda{}$ by $r$ (a margin of $< 0.01 yr^{-1}$ is not small next to $r = 0.02$). A scan over $\tau{} \in{} [0,80] yr$ cannot rule out a single-delay Hopf, because the competing demographic-delay Hopf appears only at large $\tau{}_{P} \approx{} 225 yr$ (slow $\omega{} \approx{} 0.011 yr^{-1}$) once parameters are perturbed off the knife-edge. State the location of any "no single delay destabilises" claim as a function of $\Lambda{}$/$\chi{}$, and ensure any such scan extends to $\tau{}_{P} \gtrsim{} 250 yr$ before concluding. (ii) State whether interpolation occurs for the $\tau{}_{g}/\tau{}_{p}$ (both exact multiples of $\Delta{}t$), or use a non-multiple step. (iii) Complete the scenario/parameter table (which $e$, which lags, whether $T_{b}$ is active). (iv) Analyse the trivial equilibrium $(A,P)=(0,0)$ and the no-recovery region. (v) Reconcile the "max \ensuremath{\Omega{}} not reported" footnote (\ensuremath{\Omega{}} peaks at \ensuremath{\approx{}}5.0 on $D$, \ensuremath{\approx{}}4.1 on $E$, clipped at 8.5).
\item \textbf{Characteristic equation, crossing curves, and full spectrum.} We apply the exact crossing-curve formalism (Hale \& Huang 1993; Gu, Niculescu \& Chen 2005) and a full-spectrum computation to the constant-parameter subsystem. Linearising about an interior point $A*$ of the equilibrium family $P=B(A)/e$ (defined for full harvest, $\sigma{} = 1$) gives $D(s;\tau{}_{g},\tau{}_{p})=(s-a_{1}e^{-s\tau{}_{g}}-a_{3})(s+r)-a_{E} a_{4}e^{-s\tau{}_{p}}=0$, which has a zero eigenvalue $D(0)=0$ on the whole family (a neutral continuum — no isolated interior attractor) and a positive real leading eigenvalue ($Re \lambda{} \approx{} +0.62$ at $A*=0.8$) for all delays, guaranteed by the sufficient condition $S = G\prime{}(A*) + b/b_{G} > r$ ($0.6083 > 0.02$, and $\geq{} 0.57 > r$ at $A*=0.8$). Scanning $s=i\omega{}$ yields no imaginary-axis crossing, and the zero-delay $a_{11}<r$ condition is violated everywhere ($a_{11}=G'(A*)+b/b_{G} \geq{} 0.57 > r$). The $\chi{}$ two-gain Hopf classification (derived for the interior attractor of the gross-harvest formulation) does not transfer; the crossing-curve method should be applied to the boundary equilibrium, not an interior point (\texttt{scans/r2\_char\_spectrum.png}, \texttt{scans/r2\_a11\_vs\_delay.png}).
\end{itemize}

\textbf{Comparison of formulations.} Two distinct formulations are frequently conflated in the surplus-production / sustainability literature, and they differ qualitatively. In the \textbf{gross-harvest} (Schaefer-type) formulation, total demand is removed directly from the stock; its constant-parameter subsystem has a unique interior attractor, can be stable for a single delay, and can exhibit a two-delay Hopf ($\chi{}=1$, $\tau{}_{g}*\approx{}85$, $\tau{}_{P}*\approx{}225$, $s=\pi{}/\omega{}$). In the \textbf{deficit-driven} formulation of Eq. (4), only the shortfall above the flow yield liquidates stock; its constant-parameter subsystem has a one-parameter equilibrium family $P = B(A)/e$ (no isolated interior attractor), a monotone positive-real eigenvalue (\ensuremath{\approx{}} +0.62) for every delay, no imaginary-axis crossing, and a recover/collapse basin dichotomy. The two are compared directly:

\begin{table}[htbp]
\centering
\small
\begin{tabularx}{\textwidth}{l>{\raggedright\arraybackslash}X>{\raggedright\arraybackslash}X}
\toprule
 & Gross-harvest (surplus-production) formulation & Deficit-driven formulation (Eq. 4) \\
\midrule
Equilibrium & unique interior attractor & one-parameter family $P = B(A)/e$ (no isolated attractor), under full harvest $\sigma{} = 1$ \\
Linearisation & interior point; $a_{11} < r$ can hold & interior point: $a_{11} = G'(A*) + b/b_{G} > r$ (the sufficient condition $S > r$) \\
Stability & stable for a single delay & monotone positive-real eigenvalue (\ensuremath{\approx{}} +0.62) for every delay (under $S > r$), $\sigma{} = 1$ \\
Hopf / onset & two-delay Hopf: $\chi{}=1$, $\tau{}_{g}*\approx{}85$, $\tau{}_{P}*\approx{}225$, $s=\pi{}/\omega{}$ & none: $D(0)=0$ (neutral continuum); no imaginary-axis crossing \\
Classification & $\chi{}$ two-gain Hopf (which-lag) & structural vicious cycle, no $\chi{}$; crossing-curve \ensuremath{\to{}} boundary eq. \\
Basin & stable fraction 0.506 \ensuremath{\to{}} 0.042 & recover/collapse dichotomy 39.9 \% \ensuremath{\to{}} 5.3 \% \\
\bottomrule
\end{tabularx}
\end{table}

\textbf{The "balanced-with-lags" case.} In the gross-harvest formulation the balanced case ($e = r_{opt}$) with lags oscillates (the \ensuremath{\chi{}}=1 Hopf, $\tau{}_{g}* \approx{} 85 yr$) rather than collapsing. On the deficit-driven subsystem the balanced case still lies on the continuum and its interior points are monotonically unstable, so the "no collapse at balance but oscillation with lags" dichotomy does not apply. A full-$(7)$ recompute is required before asserting either outcome. \textbf{Empirical grounding of $\tau{}_{g}$ and the delay-response sweep.}

\begin{itemize}
\item \textbf{Operational definition.} $\tau{}_{g}$ is the regeneration/recruitment-response lag — the delay between a change in the environment state $A$ and the regeneration/recruitment response that rebuilds carrying capacity ($A(t - \tau{}_{g})$). In the model's framing, "a tree takes $\tau{}_{g} \approx{} 20–80 yr$ to bear fruit." Concretely: the time from environmental change until the regeneration/recruitment response begins to restore carrying capacity, approximated by the time to reach \ensuremath{\approx{}}50 \% of pre-disturbance productive capacity. The \ensuremath{\approx{}}50 \% threshold excludes the early "green shoots" transient and matches the time to meaningful biomass/stock recovery in the cited field studies (forest AGB recovered \ensuremath{\approx{}}122 Mg ha⁻¹ within 20 yr — of order half of old-growth at many sites — median \ensuremath{\approx{}}66 yr to 90 \%; soil SOC at a new equilibrium after \~{}23 yr and \~{}17 yr). It is a conservative definitional threshold, not a fitted value; the collapse cliff sits at $\tau{}_{g} \gtrsim{} 20 yr$ either way.
\item \textbf{Field-derived band (literature banding, cited; not a full digitisation).} Forests: AGB recovery is substantial by 20 yr (\ensuremath{\approx{}}122 Mg ha⁻¹, i.e. of order half of old-growth at many sites), median \ensuremath{\approx{}}66 yr to 90 \% of old-growth (Poorter et al. 2016). Soils: SOC build-up begins \~{}5 yr and reaches a new equilibrium after \~{}23 yr (deforestation) and \~{}17 yr (grassland\ensuremath{\to{}}cropland) (Poeplau et al. 2011; IPCC 2006). Fisheries: rebuilding is slow and often incomplete — only \~{}29 \% of collapsed stocks recovered to 50 \% within 5–15 yr, most showed little change by 15 yr, and recovery is generally within \ensuremath{\approx{}}20 yr once fishing pressure is reduced to FMSY (Hutchings \& Reynolds 2004; Neubauer et al. 2013). Core band \ensuremath{\approx{}} 10–40 yr; extended \ensuremath{\approx{}} 5–60 yr; baseline $\tau{}_{g} = 30$ sits inside the band. Note the earlier readings "forest \ensuremath{\approx{}}50 \% by 20 yr" and "fisheries \ensuremath{\approx{}}12–13 yr to 50 \%" are not what these sources report; see the corrected, clearly-labelled worked conversion in SI §S5.1 ("Illustrative empirical anchoring of the regeneration timescale"), which treats the published timescales as an \emph{effective} regeneration timescale via $\tau{}_{g} \approx{} 1.44 \cdot{} t_{50}$ (\ensuremath{\approx{}}25–33 yr, representative \ensuremath{\approx{}}29 yr) and is explicitly an illustration, not a formal calibration.
\item \textbf{Delay-response sweep (baseline $\tau{}_{p}=25$).} The monotone structural instability (no Hopf, $Re \lambda{} \approx{} +0.62$) holds for every $\tau{}_{g}$; the collapse basin is robust (\ensuremath{\approx{}}5 \%) across the whole field-supported band. A 1-yr-step fine sweep over $\tau{}_{g} \in{} [15,25]$ resolves the transition: the recover fraction descends within a narrow band — 0.399 ($\tau{}_{g} \lesssim{} 17$), 0.394 (18), 0.240 (19), 0.053 ($\tau{}_{g} \gtrsim{} 20$) — so the "cliff" is a steep transition band \ensuremath{\approx{}}18–20 yr (half-way \ensuremath{\approx{}}0.24 at $\tau{}_{g} = 19$), not a single hard edge.
\end{itemize}

\begin{table}[htbp]
\centering
\small
\begin{tabularx}{\textwidth}{>{\raggedright\arraybackslash}Xllll}
\toprule
$\tau{}_{g}$ (yr) & 0–17 & 18 & 19 & 20–60 \\
\midrule
leading Re \ensuremath{\lambda{}} & 0.608–0.625 & 0.625 & 0.625 & 0.625 (constant) \\
recover fraction & 0.399 & 0.394 & 0.240 & 0.0529 \\
\bottomrule
\end{tabularx}
\end{table}

The collapse result holds for $\tau{}_{g} \gtrsim{} 20$ — the range the field data support for forests, soils, and many (non-recovered) fisheries — while the fast-regeneration regime ($\tau{}_{g} \lesssim{} 18$) returns the recovery outcome. It is not merely a one-time basin loss: §Discussion shows the recovery is not self-sustaining either. State the interval scoped, never as a flat "10–40 yr": the relevant range for the collapse mechanism is $\tau{}_{g} \gtrsim{} 20 yr$. The field-supported band lies entirely in the collapse regime for the baseline parameter set.

\begin{itemize}
\item \textbf{$\tau{}_{p}$ handling and robustness.} $\tau{}_{p}$ is held at 25 yr for the sweep; the $(\tau{}_{g}, \tau{}_{p})$ grid spans $\tau{}_{g} \in{} {10,15,18,19,20,30}$ \ensuremath{\times{}} $\tau{}_{p} \in{} {10,20,25,30,40}$. The cliff is governed by $\tau{}_{g}$: the recover fraction is 0.399 at $\tau{}_{g} = 10–15$ and 0.053 at $\tau{}_{g} = 30$ for every $\tau{}_{p}$, and at the transition the intermediate values move only slightly with $\tau{}_{p}$ (0.399–0.394 at $\tau{}_{g}=18$; 0.245–0.269 at $\tau{}_{g}=19$; 0.053–0.058 at $\tau{}_{g}=20$). So $\tau{}_{p}$ shifts the band by <1 yr across a 4\ensuremath{\times{}} range (10–40 yr) — the cliff is not a $\tau{}_{p}$ artefact. A full two-delay sweep is left to future work, but the coarse grid is no longer the only support: the flatness holds out to $\tau{}_{p} = 10$ and $40$.
\item \textbf{The cliff is $\tau{}_{p}$-independent across the whole tested set, and there is no Hopf anywhere in the delay plane.} Extending the sweep to $\tau{}_{p} \in{} {0,5,10,15,20,25,40,60}$ with $\tau{}_{g}$ at 1-yr resolution: the last recovering $\tau{}_{g}$ is 18 for every $\tau{}_{p}$ (no dimensionless $\tau{}_{g}/\tau{}_{p}$ tuning rescues the cliff), and a full $(\tau{}_{g}, \tau{}_{p})$-plane scan of the characteristic roots finds zero imaginary-axis crossings with $Re \lambda{} \approx{} +0.625$ constant (\texttt{scans/topdown\_delay\_boundary.png}). The "basin-boundary crisis, not a fold" statement is thereby confirmed over the whole plane, not just a single-delay slice. The same computation supplies the §4.5 macro-ratio results (onset $R_{B}=1.000$/$R_{A}=1.01–1.06$; silent-collapse fraction; separator balanced accuracy; rescue set), which corroborate the neutral-continuum / no-Hopf finding from the monitoring view.
\item \textbf{Conditional on baseline parameters.} The transition and recovery/collapse dichotomy are reported at the baseline $(b_{G}=0.8, A_{ref}=0.8, \rho{}=0.05, A_{max}=1.2, e=0.55, r=0.02)$ and the documented IC grid (see the scenario and parameter tables in §S4). A full global sensitivity analysis over all parameters is not performed; these are conditional statements holding while the other parameters stay at baseline, and the robustness that has been checked is recorded in §S5.
\item \textbf{Calibration outlook.} These are qualitative/regime claims from a literature-banded sweep, not a fitted forecast. A formal calibration (estimating $\tau{}_{g}$, and possibly $b$, from a small set of well-documented case studies, or a full ML/Bayesian fit) is feasible but not required for the conceptual claims and may not sharpen the conclusions given weak delay identifiability; it is left to future work.
\item \textbf{Honesty caveat.} This is a literature-banded model sweep, not a full calibration: the per-study curve digitisation is still pending, and $\tau{}_{g}$ is a lumped regeneration-response delay across heterogeneous ecosystems. The defensible claim is that the collapse mechanism operates for any $\tau{}_{g}$ in the observed band — not that $\tau{}_{g}$ equals a single measured value. Report $\tau{}_{g}$ as an interval over the observed band, not as a single number; and when a point margin is small relative to the spread of that band, do not present it as a precision result.
\end{itemize}

\section{Policy extensions (Half-Earth \& reservation)}

\begin{itemize}
\item Cap the human-available flow at $\sigma{} B$; compute both $K$ and the additional debt from that cap; report $\Omega{}$ against the allocated share (not full $B$); state whether lags/$T_{b}$ are on. This removes the impossible "$\Omega{} = 0.575$, $D = 0$" pair.
\item Corrected cap: $K = 0.5 B / e$ (not $0.5 B / r_{opt}$); verified to settle at $\Omega{} \to{} 0.500$, $P = 0.217$, and a reservation policy that keeps $E$ inside the allocated flow succeeds (no liquidation).
\item \textbf{The stabilising institutional sign is the protective one.} The vicious cycle is the mobilising / extractive sign (demand responds to a perceived shortfall by raising pressure on the base); the reservation / Half-Earth cap is the protective sign (demand is held to a cap). The design principle is therefore a sign statement, not a delay statement: the stabilising institutional response is the protective / effort-reducing one, which is why the reservation (protective) policy is the candidate to pursue.
\item The debt-repayment dichotomy is resolved: report the qualitative change (collapse \ensuremath{\to{}} oscillation for $\eta{} \gtrsim{} 0.02$) and treat $\eta{} = 0$ as the irreversible-debt benchmark.
\item \textbf{The reservation result is a nominal-path consequence.} The "$E$ inside the allocated flow succeeds, no liquidation" result is asserted under the nominal path. To make it defensible it should be scored as a robust-viability claim: declare a disturbance class (persistent productivity/$\rho{}$/$b$-shocks), define a preregistered retention rule (a governance module is kept only if it improves a declared protection-and-supply score), and apply an erosion conversion (convert the declared model defect into a certified-kernel erosion margin from the closed loop's contraction rate). Until that is done, state the result as a nominal-path consequence, not a robustly-viable policy claim.
\item \textbf{The reservation claim is certified only under a controller that observes the typed floor.} The result is scored against the nominal path. The reason it is only nominal is structural: the policy is certified for a controller that observes the typed stock floor $A$ (and the allocated flow), not merely the composite $B$. If only an aggregate composite is observed, no observation-based policy can guarantee the stock floor — the safe controls of the states consistent with a given observation may intersect emptily, so the failure is informational (which quantitative feature of the observation design — its coarseness / aggregation relative to the typed floor — is responsible) rather than dynamical.
\item \textbf{Architecture-substitution caveat.} The model analyses continuous delays $\tau{}_{g}, \tau{}_{p}$. If a policy review cadence is modelled as sample-and-hold / periodic review, do not treat it as "the delay equation sampled at $T_{r}$": the two are different operators, and moving between them can move or delete stability boundaries. Any future "review interval" extension must be a separate operator with its own monodromy, and continuous-$\tau{}$ conclusions must not be transferred to it.
\end{itemize}

\section{Demonstration}

\textbf{Fig. \texttt{IMPLEMENTED\_demo.png}} (generated by \texttt{demo\_unified.py}, the well-posed smooth-ramp solver) shows a representative overshoot run: the stock $A$, biocapacity $B$ and population $P$ all rise (a logistic overshoot / boom), then the system collapses to the extinction floor with debt $D$ building. Two model properties are visible and reproducible: the delayed-regeneration overshoot ($A$ briefly exceeds $A_{max}$, a Hutchinson-style transient here around a long $\tau{}_{g}$/$\tau{}_{p}$); and the debt \ensuremath{\to{}} degradation \ensuremath{\to{}} collapse route, in which lifting yield $b$ raises $K$, $P$, and $E$, so that $D$ accrues and $b$ is eventually eroded. \textbf{The productivity illusion ("$B$ rises while $A$ falls") — quantified, converged, and conditional.}

\begin{itemize}
\item \textbf{Method.} The reduced masking model ($dA/dt = G(A) - ramp(E-bA)/b_{G}$, $dD/dt = ramp(E-B) - \eta{}D$, $B = bA + b_{G} G(A)$, $b = (b_{0}+T_{b})e^{-\alpha{}D}$, softplus ramp $w=0.05$) was integrated with converged RK4; a 7,128-case scan plus a deficit sweep. The softplus width $w = 0.05$ is stated, and insensitivity to $w$ is reported; a mask is scored as a contiguous span with $dB/dt > 0$ and $dA/dt < 0$, converged over $dt = 0.25 \to{} 0.01$.
\item \textbf{Result.} A genuine mask exists, but only for a small initial deficit. For the favourable configuration ($\rho{}=0.05, b_{G}=0.8, b_{0}=0.5, \eta{}=0.05, \alpha{}=0.03, \kappa{}=0.2, t_{wave}=15, \Delta{}b=1.5$, $A_{0}=1.0$), the window is 5.4 yr wide at deficit $E - b_{0}A_{0} = 0.06$: across it ($t \approx{} 1.4 \to{} 6.8$ yr) $B$ rises $0.576 \to{} 0.647$ (peak) while $A$ falls $0.959 \to{} 0.858$; over the run from $t=0$ (where $A=1.00$, $B=0.578$) the peak-$B$ rise is 0.069 and $A$ falls by 0.142 (fig. \texttt{IMPLEMENTED\_demo\_masking.png}, panel a). It is converged (identical at $dt = 0.1/0.05/0.02/0.01$), so it is not an Euler artifact.
\item \textbf{The mask is bounded and deficit-limited.} Window width grows with deficit up to \ensuremath{\approx{}}5.4 yr at deficit 0.06, then collapses to zero at deficit \ensuremath{\approx{}}0.075 (fig. panel b): beyond a small overshoot the stock-liquidation feedback dominates and technology cannot lift $B$ — the run goes straight to the extinction floor with no mask.
\item \textbf{Scope of the claim.} The earlier "rising biocapacity masks a falling stock" narrative is conditional: it holds near balanced conditions, and fails under a genuine overshoot (deficit \ensuremath{\gtrsim{}}0.075, \ensuremath{\approx{}}15 \% of the initial flow yield $b_{0}A_{0}$) — precisely the Jevons-rebound logic of §5 (technology raises $K$, $P$, $E$, so debt grows and closes the window). The illusion requires the wave to outpace the debt build-up — it must arrive while the stock is still high and the cumulative deficit is still small. A slow/late technology wave merely props up a falling yield.
\end{itemize}

\textbf{The illusion is an instance of the compensatory-aggregation gap.} The weighted composite ($B$, the weak-sustainability index) satisfies its aggregate floor while the typed floor — the stock $A$ — is violated. Quantitatively it is the substitution buffer: because $B = bA + b_{G} G(A)$, the aggregate floor on $B$ can be met while the typed floor on $A$ is violated by exactly the capital growth $b_{G} G(A)$ — the same separation that appears in §4.5 as the gap between $R_{B}$ and $R_{A}$. To say it in the model's notation: the aggregate that an index reads is a weighted sum over several floors, so a deficit in one (the stock) can be masked by a surplus in another (the yield), and a transition can satisfy the aggregate the whole way along while it never satisfies the individual floors. Two consequences follow:

\begin{itemize}
\item The mask is a structural failure of any aggregate that compensates a falling base with a rising yield, not an accident of one parameter set.
\item Under exact-tube semantics (a transition is safe only if every state along it, not merely the endpoint, satisfies the constraints), the intermediate passage where $A$ falls is itself a violation even if the endpoint recovers. That is a reason to report the through-time trajectory and the boundary curve rather than only an endpoint attractor.
\item \textbf{Why Scenario E shows no illusion.} Biocapacity reaches \ensuremath{\approx{}}0.5588 in the original scenario E, but because $A$ is rising (the logistic overshoot off the low starting population), not because of the masking mechanism. So "no shown scenario exhibits the illusion" is the correct reading there.
\item \textbf{Presentation rule.} Present the illusion as narrow, deficit-bounded, and transient (fig. \texttt{IMPLEMENTED\_demo\_masking.png}), with the computed $t_{peak}$ and the critical deficit stated as numbers.
\end{itemize}

\section{Model scope}

This is a conceptual / stylised model with representative calibration — not a forecast. Its purpose is to make explicit the logic connecting accounting, deficit, lags, and the stock/yield decomposition, and to state which outcomes are emergent (falsifiable) versus imposed (logistic regeneration, bounded technology, the lags, the accounting identity). The 1961–2022 sentence is a qualified observation: the accounts are conservative (biocapacity likely overstated, overshoot understated); the series is consistent with the illusion reading but does not demonstrate it. \textbf{National Footprint Accounts (NFA) data limitation.} The GFN accounts are account-based — measured yields and conversion factors, no soil erosion, deforestation, groundwater depletion, or other cumulative degradation. Consequently the estimated biocapacity most likely overstates the real resource base and the Footprint most likely understates the true overshoot — the actual overshoot is likely larger than the accounts document. This is the strongest reason to treat the 1961–2022 series as "consistent with" the illusion reading rather than as evidence for it. The empirical programme the Discussion promises is the $d ln B = d ln b + d ln A$ decomposition with independent $A$-proxies (land cover, soil carbon, NPP). \textbf{Information-layer limit.} The signature of the illusion ("$B$ rises while $A$ falls") is contemporaneously observable (a read-off of two measured series). But whether that rise is a genuine recovery or a technology-driven mask is not identifiable at the observation time: identifying the mask requires knowing the technology/$b$-channel contribution, which is not available contemporaneously without independent proxies. Two consequences: (i) a real-time observation of $B$ rising while $A$ falls is not evidence for the mask mechanism over the recovery mechanism — both reproduce the same signature; and (ii) the empirical programme must target the $b$-channel decomposition (independent $A$-proxies), not the composite $B$ series. \textbf{The empirical programme is a flux reconstruction, not a curve-fit.} The $b$-channel (the technology contribution to yield) is an unobserved internal flux: it is not separately measured, but it is induced by the observed stock changes. The discipline is to reconstruct unobserved internal fluxes from observed stock changes and to state which fluxes are identifiable (a flux that is not a read-out of a measured stock is not recoverable). So the programme is bounded: the $b$-channel is recoverable only relative to a declared observation operator, which is why independent $A$-proxies are needed rather than a fit to the composite $B$. ---

\section{Discussion}

\subsection{Structural properties and negative results}

Several phenomena common in ecological models are absent here, and the absence is a structural property, not an oversight. Each is stated so the reader does not expect it:

\begin{table}[htbp]
\centering
\small
\begin{tabularx}{\textwidth}{>{\raggedright\arraybackslash}X>{\raggedright\arraybackslash}X>{\raggedright\arraybackslash}X>{\raggedright\arraybackslash}X}
\toprule
Phenomenon looked for & Present? (verified) & Why absent / what it is & Implication \\
\midrule
Critical slowing down (CSD) / early warning before the $\tau{}_{g}$ cliff & No & the $\tau{}_{g}$ transition is a first-order basin-boundary crisis: $Re \lambda{} \approx{} +0.62$ is constant (no eigenvalue crossing zero), and the $P$-relaxation return time is flat \ensuremath{\approx{}}150 yr for $\tau{}_{g}=0–18$ then vanishes abruptly at $\tau{}_{g}\approx{}19–20$ & do not expect a smooth warning signal before a long-$\tau{}_{g}$ collapse; monitoring a single trajectory's relaxation time will not warn \\
CSD before the fixed-liability $E$-fold & Yes, but only capital-dominated ($b_{G}\rho{}>b$) & there $B'(A*)=0$ exactly (a genuine fold, two opposite-sign-$B'$ fixed points merging); absent at baseline where $B(A)$ is monotone & Scheffer's "loss of resilience precedes the switch" applies only to this regime-scoped fold \\
Hysteresis / path dependence in the threshold & No & $A_{c}(E)$ is single-valued; for a fixed $E$ there is at most one deficit-region fixed point (no bistability), so the collapse–recovery threshold is the same going up as coming down & no path memory in the threshold; recovery is immediate once the liability is lowered (the asymmetry is in time, not in the threshold) \\
Endogenous limit cycle (sustained boom–bust) & No & the regrow-overshoot-recollapse is a single recovery-overshoot pulse (two crossings of $0.5\cdot{}A_{max}$), consistent with the monotone (no-Hopf) instability & the overshoot is not a self-oscillation; sustained boom–bust would need an external mechanism \\
Allee-type rescue threshold (a minimum viable stock) & No & no Allee term; the collapse at $\tau{}_{g}=30$ is domain-wide — every $A_{0}$ from $0.02$ to $1.00$ (including a healthy stock $A_{0}=1.0$) collapses; the recover set is a measure-zero strip ($A_{0}=A_{max}$ recovers; $1.19$ and $1.21$ both collapse) & you cannot rescue the system by starting from a higher stock — the lag is the controlling factor; only shortening it helps \\
\bottomrule
\end{tabularx}
\end{table}

\textbf{So the collapse is robust but not predictable from a single scalar: it is driven by the regeneration lag (a finite-amplitude basin erosion), not by a fold, not by Allee low-density dynamics, and it carries no generic CSD precursor.} This is more informative than a set of positive results: it tells a manager which levers do not exist (no rescue-by-stock, no early-warning signal) and which do (shorten the lag; and, only for the capital-dominated $E$-fold, watch for slowing). The result also inverts the classical complexity–stability expectation of May (1973), who showed that increasing species number and connectance in model ecosystems does not confer stability but tends to destabilise beyond a critical complexity threshold. Here instability is generated in the opposite sense: it emerges in a deliberately minimal, low-complexity system (two state variables, a single feedback loop, and an Allee-free regeneration term), so the collapse is not a by-product of model richness but follows from the deficit mechanism and the regeneration lag alone. Model extensions suggested by the negative results — stochasticity, seasonal forcing (to test for induced cycles), or an explicit Allee term (to test for a rescue route) — are future-work directions, not changes to the present deliberately minimal, deterministic, Allee-free model. \textbf{The macro-ratio monitor fails silently in the long-$\tau{}_{g}$ regime.} The $R_{B}$/$R_{A}$ ratios of §4.5 are necessary-but-not-sufficient once the lag is too long: on the documented $A_{0}\times{}P_{0}$ grid at $\tau{}_{g} = 30$, 36.5 \% (72 of 197) of collapses begin with $R_{B} < 1$ and $R_{A} < 1$ — neither ratio warns; only shortening the regeneration lag helps. (At $\tau{}_{g} = 10$, only 1.6 \% of collapses are silent.) Ratio-based monitoring therefore cannot replace the lag: it is a leading/necessary check, not a collapse forecaster. This is prediction 8. \textbf{The rescue set collapses to a measure-zero strip.} The rescue set (initial conditions that recover) is a 39.9/39.9/39.4 \% A-span of 0.100–1.300 (13 distinct $A_{0}$) at $\tau{}_{g} = 0/10/18$, then collapses to 5.3 \% with a 1.200–1.200 (1 distinct $A_{0}$) strip at $\tau{}_{g} \geq{} 20$ — i.e. at $A_{max}$ only. Combined with the absorbing collapse (§12.2), this is the quantitative form of the "cannot rescue by starting higher" statement. \textbf{The transition is a nonlocal, not a map, bifurcation.} The asymptotic one-step map (when the delays are restored to their large-time limits) has fixed points exactly on the equilibrium family ($E = 0.15 \to{} A = 0.283$; $0.30 \to{} 0.576$; $0.50 \to{} 0.986$; boundary $A_{max} = 1.2$) and its local $Re \lambda{}_{max}$ is $+0.5917 \to{} +0.6249 \to{} +0.6250$ — constant, no flip across $\tau{}_{g} \approx{} 19$. The asymptotic map therefore reproduces the fixed points but not the $\tau{}_{g}$ cliff, confirming the transition is a nonlocal basin-boundary crisis, not a map bifurcation, and that the discrete and continuous pictures do not agree as $\tau{}_{g}$ is increased. This corroborates the no-CSD reading of the delay transition, which we distinguish from the $E$-fold of §4.2. \textbf{Separator degradation is a class-imbalance artefact unless reported by balanced accuracy.} The neutral separator $P = B(A)/e$ (i.e. $R_{B} = 1$) classifies the basin at balanced 99 \% (no-delay) \ensuremath{\to{}} 81 \% ($\tau{}_{g} = 30$); a linear functional shows 96 \% (no-delay) \ensuremath{\to{}} 50 \% ($\tau{}_{g} = 30$, chance) while its raw accuracy is 94.7 \% — equal to the majority-collapse class (197/208). Report separator/ratio accuracy by balanced accuracy, never raw accuracy on the long-$\tau{}_{g}$ basin. \textbf{Recovery overshoot is real, initial-condition-dependent, and steep.} The recovery overshoot scales steeply with the lag (a fit to the computed pulse gives $ov \approx{} 0.0047\cdot{}(\tau{}_{g}/10)^4.84$, not a single $\tau{}_{g}$ law); it is an initial-condition-dependent transient, so report it as a scoped number rather than a universal amplitude.

\subsection{Recovery dynamics}

\begin{itemize}
\item \textbf{Recovery from the extinction floor is not self-sustaining at realistic lags.} Starting from the post-collapse floor at $\tau{}_{g} \geq{} 20$, the stock regrows but overshoots $A_{max}$ and re-collapses to $A_{ext}$, rather than settling at the sustainable boundary: $A_{peak}$ = 1.31 ($\tau{}_{g}=20$), 1.45 (30), 1.60 (40), 1.88 (60), and the fraction of the run spent above $0.5\cdot{}A_{max}$ falls monotonically (0.16 \ensuremath{\to{}} 0.045 \ensuremath{\to{}} 0.033 \ensuremath{\to{}} 0.020) — so the collapse is effectively absorbing, not a transient excursion. This matches the field picture of slow, incomplete, often-failed recovery (Hutchings \& Reynolds 2004; Neubauer et al. 2013). The collapse is not merely hard to reverse, it is self-sustaining under the regeneration lag. (This is separate from the masking window and from the Allee-free fixed-liability threshold — the same delay-driven liquidation, read as a recovery-blocking mechanism.)
\item \textbf{The regeneration rate sets how long recovery takes; the regeneration lag sets whether it happens at all.} Sweeping $\rho{}$ and $\tau{}_{g}$ independently from the extinction floor, the time to 50 \% of $A_{max}$ scales roughly inversely with $\rho{}$ (t50 \ensuremath{\approx{}} 145/\ensuremath{\rho{}} yr: 203 at $\rho{}=0.03$, 142 at 0.05, 105 at 0.08, 82 at 0.12, at $\tau{}_{g}=30$), while whether it recovers at all is decided by $\tau{}_{g}$ (recovered at $\tau{}_{g} \leq{} 18$, collapsed at $\tau{}_{g} \geq{} 20$, for every $\rho{}$):
\end{itemize}

\begin{table}[htbp]
\centering
\small
\begin{tabularx}{\textwidth}{l>{\raggedright\arraybackslash}X>{\raggedright\arraybackslash}X}
\toprule
$\rho{}$ (yr⁻¹) & $\tau{}_{g}=10$ (recovers) & $\tau{}_{g}=30$ (collapses) \\
\midrule
0.03 & t50 \ensuremath{\approx{}} 162 yr & t50 \ensuremath{\approx{}} 203 yr \\
0.05 & 106 & 142 \\
0.08 & 74 & 105 \\
0.12 & 56 & 82 \\
\bottomrule
\end{tabularx}
\end{table}

A manager has two genuinely different levers: accelerate the regeneration rate (restore stock-by-stock productivity) to shorten the time to recovery, versus shorten the regeneration lag (recruit/protect the seed source) to change whether recovery happens. The two do not substitute: raising $\rho{}$ makes a doomed trajectory recover faster to 50 \% but does not keep it from collapsing under a large $\tau{}_{g}$.

\begin{itemize}
\item \textbf{Recovery overshoots the old $A_{max}$ before settling (a catch-up rebound pulse).} From the floor, a system that does recover rises past $A_{max}$ before relaxing: $A_{peak}$ = 1.21 ($\tau{}_{g}=10$), 1.25 (15), 1.28 (18) — up to +0.08 beyond $A_{max}$ (fig. \texttt{scans/eco\_recovery\_insight.png}). This is the Hutchinson-style overshoot on the recovery side, and it is why a recovering system can appear to "overshoot" briefly — a transient that must not be read as a new stable overshoot. (It is distinct from the §Demonstration collapse-regime overshoot, which ends in liquidation.)
\item \textbf{Grid-dependence of the recover fraction is bounded but real.} On the same integrator, a coarse 6\ensuremath{\times{}}8 vs fine 25\ensuremath{\times{}}31 grid at $\tau{}_{p}=0$:
\end{itemize}

\begin{table}[htbp]
\centering
\small
\begin{tabular}{llll}
\toprule
$\tau{}_{g}$ (yr) & fine (0.05 step) & coarse (0.2 step) & fine / coarse \\
\midrule
30 & 0.027 & 0.104 & 0.26 \\
40 & 0.212 & 0.208 & 1.02 \\
50 & 0.301 & 0.417 & 0.72 \\
60 & 0.321 & 0.396 & 0.81 \\
\bottomrule
\end{tabular}
\end{table}

The coarse grid overstates the recover fraction near collapse (at $\tau{}_{g}=30$: 0.104 coarse vs 0.027 fine), so publish the recover fraction on the fine mesh and report the coarse value only as a bound.

\begin{itemize}
\item \textbf{The boundary-curve result is more interpretable than a percentage.} Under no delay the separatrix in $(A_{0},P_{0})$ is exactly the equilibrium family line $P_{0} = B(A_{0})/e$: "recover \ensuremath{\iff{}} $P_{0} < B(A_{0})/e$" matches the classification on 99.0 \% of the 208-cell grid (the only exceptions are $A_{0} > A_{max}$, outside the stock domain). This is the numerical witness to the neutral continuum: the recover basin is precisely the stable side of the equilibrium line. Under baseline delays $(30,25)$ the basin collapses to a thin strip $A_{0} \approx{} A_{max}$ (1.2) with $P_{0} \lesssim{} B(A_{max})/e \approx{} 1.09$, and nothing else recovers — the delay strips away the protected below-the-line region except right at the sustainable boundary (5.3 \%). Present the boundary curve rather than only the fraction.
\item \textbf{The delay-response non-monotonicity at $\tau{}_{g} \approx{} 40–60$ is real, not a grid artifact.} The recover fraction has a deep minimum at $\tau{}_{g}=30$ (\ensuremath{\approx{}}0.03) and re-opens to \ensuremath{\approx{}}0.21 (40), \ensuremath{\approx{}}0.30 (50), \ensuremath{\approx{}}0.32 (60); the coarse grid shows the same min\ensuremath{\to{}}rise shape. A boundary (separatrix) curve as a function of $(\tau{}_{g},\tau{}_{p})$ is the natural follow-up figure.
\end{itemize}

\subsection{Robustness of the main results}

\begin{itemize}
\item \textbf{Recover–collapse is robust to $b_{G}$.} On the documented grid (computed, $(30,25)$ baseline):
\end{itemize}

\begin{table}[htbp]
\centering
\small
\begin{tabularx}{\textwidth}{ll>{\raggedright\arraybackslash}X>{\raggedright\arraybackslash}X}
\toprule
$b_{G}$ & recover (no delay) & recover (baseline $(30,25)$) & \ensuremath{\Delta{}} \\
\midrule
0.4 & 0.394 & 0.053 & 0.341 \\
0.6 & 0.399 & 0.053 & 0.346 \\
0.8 & 0.399 & 0.053 & 0.346 \\
1.0 & 0.404 & 0.053 & 0.351 \\
1.2 & 0.404 & 0.053 & 0.351 \\
\bottomrule
\end{tabularx}
\end{table}

Over a 3\ensuremath{\times{}} range of the standing-stock value, the no-delay recover fraction varies only 0.394\ensuremath{\to{}}0.404 (\ensuremath{\pm{}}1.3 \%) and the baseline-delay recover fraction is pinned at 0.053 — the qualitative result is robust, not a $b_{G}$ artefact.

\begin{itemize}
\item \textbf{The monotone instability is structural} (independent of $\rho{}$). The full leading eigenvalue is essentially independent of $\rho{}$ (\ensuremath{\approx{}} +0.62, dominated by the depletion gain $a_{3} = b_{0}/b_{G}$) — itself evidence that the monotone instability is structural, not a fast–slow artifact, so the stability classification of §4.3 does not carry over.
\item \textbf{The fast–slow reduction is valid only for $\rho{} \gg{} r$ and we use the full equation.} The classification reduces the 2-D two-delay system to a scalar two-gain delayed logistic under the fast–slow separation $\rho{} \gg{} r$ (and while $a_{11} \leq{} 0$). That separation is not satisfied at realistic $\rho{}$. The table (computed at the sensitivity configuration $\tau{}_{g}=30, \tau{}_{p}=25, A*=0.8, b_{0}=0.5, b_{G}=0.8, e=0.55, r=0.02$) reports the leading real eigenvalue of the full transcendental $D(s)=0$ alongside the fast–slow $a_{11}$ and the timescale ratio $\rho{}/r$:
\end{itemize}

\begin{table}[htbp]
\centering
\small
\begin{tabularx}{\textwidth}{l>{\raggedright\arraybackslash}X>{\raggedright\arraybackslash}Xll}
\toprule
$\rho{}$ (yr⁻¹) & full-$D(s)=0$ leading $Re \lambda{}$ & fast–slow $a_{11} = G\prime{}(A*)+b/b_{G}$ & $a_{11} > r$? & $\rho{}/r$ \\
\midrule
0.02 & +0.625 & 0.618 & yes & 1.0 \\
0.05 & +0.625 & 0.608 & yes & 2.5 \\
0.10 & +0.625 & 0.592 & yes & 5.0 \\
0.50 & +0.625 & 0.458 & yes & 25 \\
1.50 & +0.625 & 0.125 & yes & 75 \\
\bottomrule
\end{tabularx}
\end{table}

Only at $\rho{} = 1.5$ ($\rho{}/r = 75$) is the fast–slow scaling justified, and even there $a_{11}$ stays positive — no qualitative flip. At realistic $\rho{} \in{} [0.02, 0.1]$ ($\rho{}/r = 1–5$) the reduction is out of its regime, so the full $D(s)=0$ must be used. We always use the full equation.

\begin{itemize}
\item \textbf{The masking band} is bounded by the deficit limits of §Demonstration (\ensuremath{\approx{}}5.4 yr at deficit 0.06; vanishes at deficit \ensuremath{\approx{}}0.075). $\alpha{}$ and $\tau{}_{p}$ were not independently swept (each needs its own grid); state $\alpha{} = 0.03$ and $\tau{}_{p} = 25$ as illustration values needing calibration, and treat the "dangerous band $(2f-1)\nu{} > 1$" as an estimated indicator.
\end{itemize}

\section{Limitations}

\begin{itemize}
\item \textbf{The sustainable point is a boundary equilibrium and stability is one-sided.} The exact $[\cdot{}]_{+}$ is non-smooth at $A = A_{max}$ (and at $E = bA$), so the "sustainable" state is a boundary of the deficit regime, not an interior point. Linearise one-sided (from the deficit/lower side only) and report that stability as such. Because the exact switch yields no imaginary-axis crossing, the non-smoothness does not create spurious oscillations. Where a smooth ramp replaces $[\cdot{}]_{+}$ (reference implementation only), fix its width $w$, state it, and report insensitivity to $w$.
\item \textbf{The masking illusion is small-deficit only, not generic.} The window maxes at \ensuremath{\approx{}}5.4 yr (deficit 0.06) and vanishes at deficit \ensuremath{\approx{}}0.075 (\ensuremath{\approx{}}15 \% of the initial flow yield $b_{0}A_{0}$), and it is converged (RK4). It does not appear for deficit \ensuremath{\gtrsim{}}0.075 and must not be cited as a general phenomenon — only as the narrow, transient, small-overshoot band of §Demonstration.
\item \textbf{Parameters are representative, not estimated; the two headline results have been sensitivity-checked.} The results are illustrative and need calibration; $b_{G}$, $\alpha{}$ and the lags are lumped, not measured. $\tau{}_{g}$ is the partial exception: it is field-banded as a regeneration/recruitment-response lag (\ensuremath{\approx{}}time to \~{}50 \% of pre-disturbance productive capacity), and the fine sweep shows the collapse basin and monotone no-Hopf instability are robust across that band for $\tau{}_{g} \gtrsim{} 20 yr$ (transition band \ensuremath{\approx{}}18–20 yr), with the extended $\tau{}_{p}$ grid not moving the cliff; the fast-regeneration regime $\tau{}_{g} \lesssim{} 18$ returns the recovery outcome. \textbf{Fit-defect disclosure.} Any parameter pinned at an optimisation bound (a fit artefact rather than an estimate) is declared as such, never presented silently as a calibration. \textbf{Interval discipline.} A point-margin small next to an uncertainty scale is reported as a coin-flip, not as skill/stability; the same wording discipline applies to any tight margin here.
\item \textbf{$\eta{}$ is set physical (non-zero); the singular $\eta{} \to{} 0$ case concerns the full model.} The constant- parameter subsystem S0 drops $D$ and $\eta{}$, so the \ensuremath{\eta{}}-singularity caveat does not apply to it. $\eta{} \to{} 0$ applies to the full $(7)$ equation, where $\eta{}$ decides whether an equilibrium exists. We set $\eta{} = 0.05 yr^{-1}$, justified as a minimal environmental regeneration / degradation-removal rate. Outcomes are insensitive to $\eta{}$ over a modest range provided $\eta{} > 0$; a reader setting $\eta{} = 0$ deliberately removes the regeneration channel and must be told it removes the equilibrium rather than perturbing it.
\end{itemize}

\section{Conclusions}

The model yields a sharp, policy-relevant conclusion. The balance point is the biocapacity ratio $R_{B} = 1$ (footprint = total biocapacity), which coincides with the neutral equilibrium family $P = B(A)/e$ and, in the capital-dominated regime, with the interior-MSY fold threshold. The commonly monitored flow-yield ratio $R_{A} = 1$ is only a leading, non-causal signal, and its lead over $R_{B}$ is governed by the flow share $\psi{}$ — so a sustainable system can legitimately run at $R_{A} > 1$. The more consequential result is negative. Once the regeneration lag is sufficiently long, neither ratio forecasts collapse: the transition is a nonlocal basin-boundary crisis carrying no critical-slowing-down precursor, the recover basin collapses to a measure-zero strip, and recovery from the extinction floor is not self-sustaining. The controlling variable is the regeneration lag, not the ratio; the manager's lever is to shorten the lag, and no early-warning signal or rescue-by-stock exists. Only in the capital-dominated $E$-fold does the classical "loss of resilience precedes the switch" structure apply, and even there it is regime-scoped. The productivity illusion — "biocapacity rises while the stock falls" — is real but narrow, deficit-bounded, and transient, and it is a structural consequence of any aggregate that compensates a falling typed floor with a rising yield. It is not identifiable contemporaneously without independent proxies for the technology channel. The paper therefore contributes a monitoring boundary ($R_{B} = 1$), a leading-indicator caution ($R_{A}$), and a precise statement of when monitoring fails (long regeneration lag), together with the policy-relevant levers (shorten the lag; adopt the protective institutional sign) and the levers that do not exist (no rescue-by-stock, no early warning). ---

\section{References}

\begin{itemize}
\item Anderies, J. M. (2000). On modeling human behavior and institutions in simple ecological economic systems. \emph{Ecological Economics}, 35(3), 393–412.
\item Blomqvist, L., Brook, B. W., Ellis, E. C., Kareiva, P. M., Nordhaus, T. \& Shellenberger, M. (2013). Does the shoe fit? Real versus imagined ecological footprints. \emph{PLoS Biology}, 11(11), e1001700.
\item Borucke, M., Moore, D., Cranston, G., Gracey, K., Iha, K., Larson, J., Lazarus, E., Morales, J. C., Wackernagel, M. \& Galli, A. (2013). Accounting for demand and supply of the biosphere's regenerative capacity: The National Footprint Accounts' underlying methodology and framework. \emph{Ecological Indicators}, 24, 518–533.
\item Brander, J. A. \& Taylor, M. S. (1998). The simple economics of Easter Island: A Ricardo–Malthus model of renewable resource use. \emph{American Economic Review}, 88(1), 119–138.
\item Cohen, J. E. (1995). Population growth and Earth's human carrying capacity. \emph{Science}, 269(5222), 341–346.
\item Dasgupta, P. (2021). \emph{The Economics of Biodiversity: The Dasgupta Review.} London: HM Treasury.
\item Diamond, J. (2005). \emph{Collapse: How Societies Choose to Fail or Survive.} New York: Viking.
\item Galli, A., Giampietro, M., Goldfinger, S., Lazarus, E., Lin, D., Saltelli, A., Wackernagel, M. \& Müller, F. (2016). Questioning the ecological footprint. \emph{Ecological Indicators}, 69, 224–232.
\item Giampietro, M. \& Saltelli, A. (2014). Footprints to nowhere. \emph{Ecological Indicators}, 46, 610–621.
\item Gu, K., Niculescu, S.-I. \& Chen, J. (2005). On stability of time-delay systems with perturbed parameters. (Crossing-curve formalism.)
\item Hale, J. K. \& Huang, W. (1993). Global geometry of the stable regions for two delay differential equations. \emph{Journal of Mathematical Analysis and Applications}, 178(2), 344–362.
\item Haberl, H. \& Aubauer, H. P. (1992). (Application of the delayed-logistic framework to human population/load dynamics.)
\item Hutchings, J. A. \& Reynolds, J. D. (2004). Marine fish population collapses: consequences for recovery and extinction risk. \emph{BioScience}, 54(4), 297–309.
\item Hutchinson, G. E. (1948). Circular causal systems in ecology. \emph{Annals of the New York Academy of Sciences}, 50(4), 221–246.
\item Kuang, Y. (1993). \emph{Delay Differential Equations, with Applications in Population Dynamics.} Academic Press.
\item Lin, D., Hanscom, L., Murthy, A., Galli, A., Evans, M., Neill, E., et al. (2018). Ecological footprint accounting for countries: updates and results of the National Footprint Accounts. \emph{Sustainability}, 10(12).
\item May, R. M. (1973). Stability and complexity in model ecosystems. \emph{Princeton University Press.}
\item Meadows, D. H., Meadows, D. L., Randers, J. \& Behrens, W. W. (1972). \emph{The Limits to Growth: A Report for the Club of Rome's Project on the Predicament of Mankind.} New York: Universe Books.
\item Meadows, D. H., Meadows, D. L., Randers, J. \& Behrens, W. W. (1974). \emph{Dynamics of Growth in a Finite World.} Cambridge, MA: Wright-Allen Press.
\item Meadows, D. H., Randers, J. \& Meadows, D. L. (2004). \emph{The Limits to Growth: The 30-Year Update.} White River Junction, VT: Chelsea Green.
\item Meyer, P. S. (1994). Bi-logistic growth. \emph{Technological Forecasting and Social Change}, 47(1), 89–102.
\item Meyer, P. S. \& Ausubel, J. H. (1999). Carrying capacity: A model with logistically varying limits. \emph{Technological Forecasting and Social Change}, 61(3), 209–214.
\item Neubauer, P., et al. (2013). Resilience of recovering fish populations. \emph{Fish and Fisheries}, 14(3).
\item Poeplau, C., et al. (2011). Temporal dynamics of soil organic carbon after land-use change. \emph{Global Change Biology}, 17(7), 2415–2427.
\item Poorter, L., et al. (2016). Biomass resilience of Neotropical secondary forests. \emph{Nature}, 530(7590), 211–214.
\item Schaefer, M. B. (1954). Some aspects of the dynamics of populations important to the management of the commercial marine fisheries. \emph{Bulletin of the Inter-American Tropical Tuna Commission}, 1(2), 25–56.
\item Scheffer, M., Carpenter, S., Foley, J. A., Folke, C. \& Walker, B. (2001). Catastrophic shifts in ecosystems. \emph{Nature}, 413(6856), 591–596.
\item Shampine, L. F. \& Thompson, S. (2001). Solving DDEs in MATLAB. \emph{Applied Numerical Mathematics}, 37(4), 441–458.
\item Simon, J. L. (1996). \emph{The Ultimate Resource 2.} Princeton, NJ: Princeton University Press.
\item van den Bergh, J. C. J. M. \& Grazi, F. (2015). Reply to the first systematic response by the Global Footprint Network to criticism: a real debate finally? \emph{Ecological Indicators}, 58, 458–463.
\item van den Bergh, J. C. J. M. \& Rietveld, P. (2004). Reconsidering the limits to world population: Meta-analysis and meta-prediction. \emph{BioScience}, 54(3), 195–204.
\item Wackernagel, M. \& Rees, W. (1996). \emph{Our Ecological Footprint: Reducing Human Impact on the Earth.} New Society Publishers.
\item Wackernagel, M., et al. (2002). Tracking the ecological overshoot of the human economy. \emph{PNAS}, 99(14).
\item Abaee, A. (2026). The Limits of Compensatory Aggregation: A Formal Separation of Weak and Strong Sustainability Assessment. Zenodo. \url{https://doi.org/10.5281/zenodo.22545740} — compensatory aggregation.
\item Abaee, A. (2026). Typed Flux Ledgers and Depletion Arithmetic: Conservation, Componentwise Diagnostics, and the Semantics of Depletion Horizons. Zenodo. \url{https://doi.org/10.5281/zenodo.22554177} — typed flux ledgers.
\item Abaee, A. (2026). An Obstruction Calculus for Viability under Incomplete Observation. Zenodo. \url{https://doi.org/10.5281/zenodo.22552616} — incomplete-observation viability.
\item Abaee, A. (2026). Delay-Induced Regime Change in Harvested Stocks: The Mobilising and Protective Channels of Institutional Feedback, and the Review Interval as Control. Zenodo. \url{https://doi.org/10.5281/zenodo.22554217} — mobilising vs. protective controller sign; review interval as control.
\item Abaee, A. (2026). Periodic Review as Sampled Governance: Sample-and-Hold Dynamics of Assessment-Driven Effort Control, a Selected 42-Stock Spectral Screen, and the Northern Cod Case. Zenodo. \url{https://doi.org/10.5281/zenodo.22554297} — periodic review / sample-and-hold governance.
\item Abaee, A. (2026). Robust viability of the 2J3KL limit reference point under a surplus-production map: policy scoring, expansion, and when catch cannot help. Zenodo. \url{https://doi.org/10.5281/zenodo.22552060} — policy scoring; negative-certificate and interval-discipline methods.
\item Abaee, A. (2026). Does a surplus-production ladder improve forecasts of Northern cod? A scored test on NAFO 2J3KL. Zenodo. \url{https://doi.org/10.5281/zenodo.22553609} — surplus-production forecast scoring.
\item Abaee, A. (2026). Governance operators and viability kernels of the Edwards Aquifer: an intervention-selection test at J-17. Zenodo. \url{https://doi.org/10.5281/zenodo.22553311} — governance operators / viability kernels.
\item Abaee, A. (2026). Does a one-pool water-balance model improve forecasts of Edwards Aquifer head? A scored test at J-17. Zenodo. \url{https://doi.org/10.5281/zenodo.22552680} — one-pool water-balance forecast scoring.
\end{itemize}

\textbf{(Data availability.)} This study is a mathematical and computational analysis and reports no new empirical dataset. All numerical results, parameter sweeps, and figures were generated by the accompanying model code run from stated initial conditions and parameters; the code (see Reproducibility below) is the primary artefact, and its generated outputs are written to \texttt{data/topdown\_results.json}. No primary observational data were collected. The only empirical inputs are published values used to bound the regeneration lag, which are drawn from the National Footprint Accounts and the field studies cited in the References and are not reproduced or re-distributed here; readers should consult those sources for the underlying data. \textbf{(Reproducibility.)} The model is implemented as \texttt{model\_sims/corrected.py}, \texttt{model\_sims/topdown.py}, \texttt{model\_sims/r1\_basin.py}, and \texttt{model\_sims/char\_eq.py} (drivers \texttt{model\_sims/\_run\_topdown.py}, \texttt{demo\_unified.py}, \texttt{mask\_rk4.py}); the results are written to \texttt{data/topdown\_results.json} and rendered as \texttt{scans/topdown\_macro\_ratios.png}, \texttt{scans/topdown\_ratio\_separation.png}, and \texttt{scans/topdown\_delay\_boundary.png}. All numeric results report their integration protocol and termination convention (§Numerics). \textbf{(Supplementary material.)} The following files accompany this manuscript and are cited where their content is used:

\begin{itemize}
\item \texttt{supplementary/SUPPLEMENTARY\_information.md} — \textbf{Supplementary Information (S1):} the full model specification, the symbol table and units (§2.1), the non-dimensionalisation and scaling, the analytic derivations (MSY, the fixed-liability threshold, the equilibrium family $P = B(A)/e$, the stability classification), the scenario and parameter tables, the sensitivity/robustness record, the prediction\ensuremath{\to{}}section map, and the supplementary references (§S1–§S7).
\item \texttt{supplementary/FIGURE\_CAPTIONS.md} — \textbf{Captions for supporting figures S1–S14}, each linked to the script that generates it.
\item \texttt{supplementary/REPRODUCTION\_GUIDE.md} — \textbf{Code map and reproduction guide:} dependencies, file/folder structure, and step-by-step commands to rebuild every numeric result and figure.
\item \texttt{supplementary/DATA\_AVAILABILITY.md} — \textbf{Data and code availability statement}, including a deposit/repository recommendation.
\item \texttt{supplementary/ABSTRACT\_submission.tex} — \textbf{The abstract as LaTeX} (mathematics kept as LaTeX for the submitted manuscript; 300 words).
\item \texttt{supplementary/FIGURES/S1–S14.png} — \textbf{Supporting figures S1–S14} (see the captions file): S1 feedback/causal-loop diagram; S2 macro-ratio plane ($R_{B}$ vs $R_{A}$); S3 flow-share separation; S4 delay-boundary / $\tau{}_{g}$ cliff; S5 representative overshoot run; S6 productivity-illusion / masking window; S7 recovery vs collapse; S8 recovery insight (rate vs lag); S9 characteristic spectrum; S10 $a_{11}$ vs delay; S11 sustainable-yield regimes; S12 basin heatmap; S13 recover-fraction vs $\tau{}_{g}$; S14 graphical abstract.
\end{itemize}

% ---------------- Declarations ----------------
\section*{Declarations}

\subsection*{Data availability}
The code and data that support the findings of this study are available in the accompanying supplementary package (Supplementary Information, supporting figures S1--S14, a reproduction guide, and a data/code availability statement). A version of this manuscript together with the supplementary package is deposited at \url{https://zenodo.org/records/22554480}. This study is a mathematical and computational analysis and reports no new empirical dataset. All numerical results, parameter sweeps, and figures were generated by the accompanying model code run from stated initial conditions and parameters; the code is the primary artefact, and its generated outputs are written to \texttt{data/topdown\_results.json}. No primary observational data were collected. The only empirical inputs are published values used to bound the regeneration lag, which are drawn from the National Footprint Accounts and the field studies cited in the References.

\subsection*{Funding}
This research received no specific grant from any funding agency in the public, commercial, or not-for-profit sectors. The author is an independent researcher and received no external funding for this work.

\subsection*{Competing interests}
The author declares that he has no known competing financial interests or personal relationships that could have appeared to influence the work reported in this paper.

\subsection*{Declaration of generative AI and AI-assisted technologies in the writing process}
During the preparation of this work the author used AI-assisted tools for editorial support, including language refinement, formatting, and computational-reproducibility checks on the accompanying model code. After using these tools, the author reviewed and validated the content as needed and takes full responsibility for the content of the publication.

\end{document}
